\documentclass[11pt]{article}

% ---------------------------------------------------------------------------
%  A Jacobian-Weighted Volumetric Density Correction for Mass Conservation
%  in Curvilinear Lattice-Boltzmann Flows
%
%  Compiles with: pdflatex (1-2 passes). No bibtex required.
% ---------------------------------------------------------------------------

\usepackage{amsmath}
\usepackage{amssymb}
\usepackage{amsthm}
\usepackage[margin=1in]{geometry}
\usepackage{booktabs}
\usepackage{listings}
\usepackage{xcolor}
\usepackage{graphicx}
\usepackage[colorlinks=true,linkcolor=blue,citecolor=blue,urlcolor=blue]{hyperref}

% --- CJK (Traditional Chinese abstract). Use CJKutf8 which ships with TeX Live
%     and works with pdflatex without external setup. ---
\usepackage{CJKutf8}

% --- Theorem environments ---
\newtheorem{proposition}{Proposition}
\newtheorem{theorem}{Theorem}
\newtheorem{lemma}{Lemma}
\newtheorem{corollary}{Corollary}
\theoremstyle{definition}
\newtheorem{definition}{Definition}
\theoremstyle{remark}
\newtheorem{remark}{Remark}

% --- Listings configuration for C++/CUDA ---
\definecolor{codegray}{gray}{0.5}
\definecolor{codekw}{rgb}{0.13,0.13,0.6}
\definecolor{codecomment}{rgb}{0.0,0.45,0.0}
\definecolor{codestr}{rgb}{0.6,0.1,0.1}
\lstset{
  language=C++,
  basicstyle=\ttfamily\footnotesize,
  keywordstyle=\color{codekw}\bfseries,
  commentstyle=\color{codecomment}\itshape,
  stringstyle=\color{codestr},
  numbers=left,
  numberstyle=\tiny\color{codegray},
  numbersep=6pt,
  breaklines=true,
  breakatwhitespace=false,
  showstringspaces=false,
  tabsize=2,
  frame=single,
  columns=fullflexible,
  captionpos=b
}

% --- Path / inline-code convenience macro ---
\providecommand{\path}[1]{\texttt{#1}}

\title{A Jacobian-Weighted Volumetric Density Correction for\\
       Mass Conservation in Curvilinear Lattice-Boltzmann Flows}

\author{%
  Corresponding author\\
  \small Periodic-Hill D3Q19 GILBM Solver Group
}
\date{\today}

\begin{document}
\maketitle

% ===========================================================================
%  ABSTRACT (English) + 摘要 (Traditional Chinese) + Keywords
% ===========================================================================
\begin{abstract}
\noindent
Body-fitted (curvilinear) lattice-Boltzmann (LBM) simulations of wall-bounded
turbulence accumulate a slow, mesh-dependent drift of the global mean density.
The drift is not a physical mass source but the residual compressibility error
of the weakly-compressible time stepping, and on a stretched grid it cannot be
removed by the naive device of summing nodal densities with unit weight: such a
sum integrates over \emph{computational} space and weights each physical region
by the inverse local Jacobian, biasing the controlled mean toward the
node-clustered (small-cell) zones. We formulate the only mesh-invariant remedy:
a \emph{volume-weighted} density governor that regulates the genuinely physical
mass $M=\int_\Omega\rho\,dV$. Each lattice node is assigned a median-dual
control volume $w_p$, and the volume-weighted mean
$\langle\rho\rangle_V=\sum_p w_p\rho_p/\sum_p w_p$ is restored to unity by a
single, spatially uniform shift injected into the rest population $f_0$. We prove
three exact properties of the corrector: the median-dual weights form an exact
partition of unity ($\sum_p w_p=V_{\mathrm{mesh}}$), the shift is
momentum-silent (the rest velocity is $\mathbf e_0=\mathbf0$), and the unit mean
is recovered in a single application (unit-gain fixed point). The cell volumes
that build $w_p$ are computed two ways: an exact straight-chord Shoelace area
(Method~0) and a $3\times3$ Gauss--Legendre quadrature of the in-plane Jacobian
$J_{2D}$ reconstructed by degree-5 Lagrange interpolation (Method~1). A
host-only verification harness confirms the Gauss--Legendre rule is exact through
degree~5, that Method~1 converges at sixth order on a curved, genuinely
two-dimensional (non-separable) test map while the chorded Method~0 is locked at
$O(h^2)$, and that the mass
governor restores $\langle\rho\rangle_V=1$ to reduction-rounding precision
($\sim\!10^{-12}$ at production scale, $\sim\!10^{-16}$ in the small verification
harness) while preserving momentum exactly. We give a complete
theory$\leftrightarrow$code map for the D3Q19 GILBM implementation and an
independent record verifying that the MRT collision operator and Guo forcing
inject no spurious mass, leaving the volumetric corrector as the sole, controlled
mass channel.
\end{abstract}

\begin{center}
\begin{CJK}{UTF8}{bsmi}
\textbf{摘要}\\[4pt]
\end{CJK}
\begin{minipage}{0.92\textwidth}
\begin{CJK}{UTF8}{bsmi}
\small
貼體（曲線）格子波茲曼（LBM）模擬牆束湍流時，全域平均密度會隨時間緩慢漂移，且此漂移與網格相關。
該漂移並非物理質量源，而是弱可壓縮時間推進殘留的可壓縮性誤差；在拉伸網格上，無法以「對節點密度取單位權重求和」這種素樸方式消除：此種求和實為在計算空間上積分，會以局部 Jacobian 的倒數加權每個物理區域，使受控平均偏向節點密集（小網格）區。本文提出唯一與網格無關的修正法：以體積加權的密度調節器，調控真正的物理質量
$M=\int_\Omega\rho\,dV$。每個格點賦予中位對偶（median-dual）控制體積 $w_p$，並透過注入靜止分布
$f_0$ 的單一均勻位移，將體積加權平均
$\langle\rho\rangle_V=\sum_p w_p\rho_p/\sum_p w_p$ 一次回復為 1。我們證明修正子的三項精確性質：中位對偶權重構成精確的單位分割（$\sum_p w_p=V_{\mathrm{mesh}}$）、位移不影響動量（靜止速度 $\mathbf e_0=\mathbf0$）、以及單次套用即回復單位平均（單位增益不動點）。建立 $w_p$ 所需的網格體積以兩法計算：精確直弦 Shoelace 面積（方法 0），以及對以五次 Lagrange 內插重建之平面內 Jacobian $J_{2D}$ 施以 $3\times3$ Gauss--Legendre 求積（方法 1）。獨立主機端驗證確認該求積對五次以下多項式精確，方法 1 在一個真正二維（不可分離）的曲線測試映射上達六階收斂、而直弦方法 0 鎖在 $O(h^2)$，且質量調節器在保持動量精確守恆下將 $\langle\rho\rangle_V$ 回復至約化捨入精度（生產規模約 $\sim\!10^{-12}$、小型驗證程式約 $\sim\!10^{-16}$）。文中並給出 D3Q19 GILBM 實作的完整理論$\leftrightarrow$程式碼對照，以及驗證 MRT 碰撞算子與 Guo 體積力不注入任何虛假質量的獨立紀錄，使體積修正子成為方案中唯一受控的質量通道。
\end{CJK}
\end{minipage}
\end{center}

\vspace{4pt}
\noindent\textbf{Keywords:} lattice Boltzmann method; mass conservation;
curvilinear / body-fitted mesh; median-dual control volume; Gauss--Legendre
quadrature; Jacobian metric; partition of unity; periodic hill; MRT collision;
Guo forcing.

\tableofcontents

% ===========================================================================
\section{Introduction}
\label{sec:intro}

The lattice-Boltzmann method (LBM) has matured into a competitive solver for
wall-bounded turbulence, particularly in the weakly-compressible,
single-relaxation- or multiple-relaxation-time (MRT) form
\cite{Qian1992,dHumieres2002,Kruger2017,Succi2001}. Its native Cartesian lattice, however,
is awkward for the smooth, curved walls that characterise canonical separated
flows such as the periodic-hill benchmark \cite{Frohlich2005,Breuer2009}. Two
classical routes restore geometric flexibility: interpolation-supplemented LBM
on body-fitted meshes \cite{He1997,FilippovaHanel1998}, and finite-volume
formulations on unstructured/curvilinear grids \cite{Barth1991}. In all of these
the populations are advanced on a logical (computational) lattice whose physical
image is stretched and curved.

A subtle consequence of this stretching is a slow, mesh-dependent drift of the
global mean density. In a weakly-compressible scheme the density is a dynamical
field, and the time stepping leaves a residual compressibility error that, on a
non-uniform body-fitted grid, accumulates into a global offset of
$\langle\rho\rangle$ away from its lattice reference value of unity. Left
unchecked, this offset perturbs the pressure (through the equation of state
$p=c_s^2\rho$) and contaminates long-time turbulence statistics.

The naive cure---summing the nodal densities with \emph{unit} weight and
shifting the field to drive that sum to a target---is \emph{wrong} on a stretched
mesh. As we show in \S\ref{sec:T1}, a unit-weight sum is a quadrature of the
density over \emph{computational} space, and therefore weights each physical
sub-region by the \emph{inverse} local Jacobian. Where grid lines cluster (near
the wall and in the lee of the hill) the metric is small and the unit-weight sum
over-counts; where cells are coarse it under-counts. Regulating such a biased
surrogate drives the density toward a spatially skewed mean and injects a
spurious, mesh-dependent body force into the momentum equation.

The correct invariant to monitor is the genuinely physical mass
$M=\int_\Omega\rho\,dV$, whose discrete form requires each node to carry a
\emph{physical control volume} $w_p$. This paper develops, analyses, verifies,
and documents the implementation of a \emph{Jacobian-weighted volumetric density
corrector} built on this principle. Its contributions are:
\begin{itemize}
\item a continuous/discrete theory (\S\ref{sec:theory}) showing that only the
  volume-weighted mean $\langle\rho\rangle_V$ is a mesh-invariant target, and
  proving that a uniform shift into the rest population $f_0$ restores
  $\langle\rho\rangle_V=1$ in one step while preserving momentum exactly;
\item a median-dual control-volume construction with an exact
  partition-of-unity identity (\S\ref{sec:weights});
\item two quadratures of the per-cell area---an exact straight-chord Shoelace
  area and a $3\times3$ Gauss--Legendre integration of the in-plane Jacobian
  reconstructed by degree-5 Lagrange interpolation---together with a full error
  budget (\S\ref{sec:quadrature}--\S\ref{sec:numerical-analysis})
  \cite{StoerBulirsch2002,DavisRabinowitz1984};
\item an independent host-only verification harness confirming GL-3 exactness,
  sixth-order convergence of the Jacobian method, and algebraically exact,
  momentum-preserving mass restoration (\S\ref{sec:experiments});
\item a faithful theory$\leftrightarrow$code map for the D3Q19 GILBM solver
  (\S\ref{sec:impl}) and a verification record for the MRT collision operator and
  Guo forcing \cite{GuoZhengShi2002} establishing that the corrector is the sole
  controlled mass channel (\S\ref{sec:mrt}, Appendix~\ref{app:mrt}).
\end{itemize}

% ===========================================================================
\section{Governing framework and notation}
\label{sec:framework}

\paragraph{Lattice and moments.}
The solver advances a D3Q19 lattice-Boltzmann scheme with $N_Q=19$ discrete
velocities $\mathbf e_q$ ($q=0,\dots,18$), rest velocity
$\mathbf e_0=(0,0,0)$, and lattice weights $W_q$ with $W_0=\tfrac13$. The
hydrodynamic moments are the velocity-space sums
\begin{equation}
\rho=\sum_{q=0}^{18} f_q,
\qquad
\rho\,\mathbf u=\sum_{q=0}^{18}\mathbf e_q\,f_q .
\label{eq:moments-master}
\end{equation}
Collision is performed by an MRT operator on the moment basis $M$ with
relaxation matrix $S$, projected to velocity space as $K=M^{-1}SM$; the body
force is applied by Guo forcing \cite{GuoZhengShi2002} at a compile-time Hermite
order \texttt{FORCE\_HERMITE\_ORDER}. The conserved (density and momentum)
moments are relaxed at rate zero; the verification of this property is recorded
in \S\ref{sec:mrt} and Appendix~\ref{app:mrt}.

\paragraph{Curvilinear map (unified convention).}
Throughout the paper we use a single coordinate convention. Let
$(\xi,\eta,\zeta)$ be the logical (computational) coordinates, identified with
the integer lattice indices $(j,i,k)$ as
\begin{center}
\begin{tabular}{@{}lll@{}}
\toprule
logical coordinate & lattice index & role \\
\midrule
$\xi$ & $j$ & streamwise \\
$\eta$ & $i$ & spanwise (homogeneous) \\
$\zeta$ & $k$ & wall-normal (carries the hill) \\
\bottomrule
\end{tabular}
\end{center}
The body-fitted mesh is an \emph{extrusion}: the spanwise physical coordinate
depends on the spanwise index alone, $x=x(\eta)$, with straight grid lines and
(non-uniform) spacing $\Delta x_i=x_{i+1}-x_i$, while the cross-section
$(y,z)=\bigl(y(\xi,\zeta),\,z(\xi,\zeta)\bigr)$ is a genuinely two-dimensional
curvilinear plane, $z$ being the wall-normal physical coordinate. The map is
\begin{equation}
\Phi:\;(\xi,\eta,\zeta)\;\longmapsto\;\bigl(x(\eta),\,y(\xi,\zeta),\,z(\xi,\zeta)\bigr).
\label{eq:map}
\end{equation}
Because $x$ depends on the spanwise variable only and $(y,z)$ on the in-plane
variables only, the $3\times3$ Jacobian is block-diagonal and its determinant
factors:
\begin{equation}
\mathcal{J}_{3}
\;=\;\det\frac{\partial(x,y,z)}{\partial(\xi,\eta,\zeta)}
\;=\;\pm\,x'(\eta)\,J_{2D},
\qquad
J_{2D}\;\equiv\;\frac{\partial(y,z)}{\partial(\xi,\zeta)}
\;=\;y_{\xi}z_{\zeta}-y_{\zeta}z_{\xi},
\label{eq:jacobian-block}
\end{equation}
where subscripts denote partial derivatives and $J_{2D}$ is the in-plane
(cross-section) Jacobian computed and stored by the solver (code symbol
\texttt{J\_2D}). The overall sign of $\mathcal{J}_3$ is orientation-dependent
(here $x'(\eta)$ occupies the second column, contributing a cofactor sign
$(-1)^{1+2}$); it is immaterial in practice because the code builds the volume
element from the absolute value (\texttt{fabs}). The physical volume element is
therefore
\begin{equation}
dV \;=\; |\mathcal{J}_{3}|\,d\xi\,d\eta\,d\zeta
      \;=\; |x'(\eta)|\,J_{2D}\,d\xi\,d\eta\,d\zeta ,
\label{eq:dV}
\end{equation}
i.e.\ a physical cell volume is the spanwise chord length $|\Delta x_i|$ times an
area measured in the $(y,z)$ cross-section. Nodes are indexed $p=(i,j,k)$ and
mesh cells $c=(i_c,j_c,k_c)$; $h$ denotes a representative physical cell size,
$n\sim10^6$ the number of interior reduction nodes, and
$\varepsilon=2.2\times10^{-16}$ the double-precision unit roundoff.

The next lemma justifies why the in-plane factor $J_{2D}$ of
\eqref{eq:jacobian-block} is the correct \emph{area density} of the cross-section:
it is precisely the square root of the determinant of the first fundamental form,
so the determinant---not the product of the per-direction scale factors---is the
quantity that the cell-volume quadratures of \S\ref{sec:quadrature} must integrate.

\begin{lemma}[Metric form of the in-plane area element]
\label{lem:area-metric}
Let the in-plane map $(\xi,\zeta)\mapsto\bigl(y(\xi,\zeta),z(\xi,\zeta)\bigr)$ of
\eqref{eq:map} have covariant basis vectors
$\mathbf r_\xi=(y_\xi,z_\xi)$ and $\mathbf r_\zeta=(y_\zeta,z_\zeta)$, with scale
factors (Lam\'e coefficients)
\begin{equation}
h_\xi=|\mathbf r_\xi|=\sqrt{y_\xi^2+z_\xi^2},
\qquad
h_\zeta=|\mathbf r_\zeta|=\sqrt{y_\zeta^2+z_\zeta^2},
\label{eq:lame}
\end{equation}
and first-fundamental-form (metric) components
$g_{\xi\xi}=h_\xi^2$, $g_{\zeta\zeta}=h_\zeta^2$, and
$g_{\xi\zeta}=\mathbf r_\xi\!\cdot\!\mathbf r_\zeta=y_\xi y_\zeta+z_\xi z_\zeta$.
Then the physical area element of the cross-section is
\begin{equation}
dA=\sqrt{\det g}\;d\xi\,d\zeta,
\qquad
\det g=g_{\xi\xi}\,g_{\zeta\zeta}-g_{\xi\zeta}^{\,2},
\label{eq:area-metric}
\end{equation}
and its area density equals the in-plane Jacobian of \eqref{eq:jacobian-block} in
absolute value,
\begin{equation}
\sqrt{\det g}
=\bigl|y_\xi z_\zeta-y_\zeta z_\xi\bigr|
=|J_{2D}|
=h_\xi h_\zeta\sin\varphi,
\label{eq:area-jac}
\end{equation}
where $\varphi$ is the angle between the coordinate-line tangents
$\mathbf r_\xi$ and $\mathbf r_\zeta$. In particular
$|J_{2D}|=h_\xi h_\zeta$ (the scale-factor product) holds if and only if the
coordinate lines are orthogonal, i.e.\ $g_{\xi\zeta}=0$, $\varphi=\pi/2$.
\end{lemma}

\begin{proof}
Expand the determinant of the metric directly,
\begin{equation}
\det g
=(y_\xi^2+z_\xi^2)(y_\zeta^2+z_\zeta^2)
 -(y_\xi y_\zeta+z_\xi z_\zeta)^2 .
\label{eq:detg-expand}
\end{equation}
Expanding,
$(y_\xi^2+z_\xi^2)(y_\zeta^2+z_\zeta^2)=y_\xi^2 y_\zeta^2+y_\xi^2 z_\zeta^2+z_\xi^2 y_\zeta^2+z_\xi^2 z_\zeta^2$
and
$(y_\xi y_\zeta+z_\xi z_\zeta)^2=y_\xi^2 y_\zeta^2+2\,y_\xi y_\zeta z_\xi z_\zeta+z_\xi^2 z_\zeta^2$.
The pure terms $y_\xi^2 y_\zeta^2$ and $z_\xi^2 z_\zeta^2$ cancel between the two,
leaving
\begin{equation}
\det g
=y_\xi^2 z_\zeta^2+z_\xi^2 y_\zeta^2-2\,y_\xi z_\zeta\,y_\zeta z_\xi
=\bigl(y_\xi z_\zeta-y_\zeta z_\xi\bigr)^2
=J_{2D}^{\,2},
\label{eq:lagrange-collapse}
\end{equation}
which is the two-dimensional Lagrange identity
$|\mathbf a|^2|\mathbf b|^2-(\mathbf a\!\cdot\!\mathbf b)^2=(\mathbf a\times\mathbf b)^2$
for $\mathbf a=\mathbf r_\xi$, $\mathbf b=\mathbf r_\zeta$. Taking the positive
square root gives $\sqrt{\det g}=|J_{2D}|$, and the standard surface-area element
$dA=\sqrt{\det g}\,d\xi\,d\zeta$ then yields \eqref{eq:area-metric}--\eqref{eq:area-jac}.
Finally, writing $\cos\varphi=g_{\xi\zeta}/(h_\xi h_\zeta)$,
\begin{equation}
\det g
=h_\xi^2 h_\zeta^2\bigl(1-\cos^2\varphi\bigr)
=\bigl(h_\xi h_\zeta\sin\varphi\bigr)^2,
\label{eq:detg-sin}
\end{equation}
so $|J_{2D}|=h_\xi h_\zeta\sin\varphi$, with equality to $h_\xi h_\zeta$ exactly
when $\sin\varphi=1$, i.e.\ $g_{\xi\zeta}=0$.
\end{proof}

\begin{remark}[Why the scale-factor product over-counts a skewed grid]
\label{rem:skew}
The naive product $h_\xi h_\zeta\,d\xi\,d\zeta$ is the area of the
\emph{rectangle} with sides $h_\xi\,d\xi$ and $h_\zeta\,d\zeta$; the actual cell
is the \emph{parallelogram} spanned by $\mathbf r_\xi\,d\xi$ and
$\mathbf r_\zeta\,d\zeta$, whose area is $h_\xi h_\zeta\sin\varphi=|J_{2D}|$ by
\eqref{eq:area-jac}. The two agree only on an orthogonal grid; on a skewed grid
($g_{\xi\zeta}\ne0$) the scale-factor product over-counts the area by the factor
$1/\sin\varphi$. The Steger--Sorenson elliptic (Poisson) generator promotes but
does not enforce orthogonality, so in general $g_{\xi\zeta}\neq0$ (the grid is only
near-orthogonal); wherever $g_{\xi\zeta}\neq0$ the determinant $J_{2D}$, and not the
product $h_\xi h_\zeta$, is the correct area density. The discrete Shoelace area
(Method~0) likewise returns the true straight-chord parallelogram/quadrilateral
area and so also accounts for skew; both cell-volume estimators of
\S\ref{sec:quadrature} therefore integrate the metric area element $|J_{2D}|$ of
Lemma~\ref{lem:area-metric}, never the scale-factor product.
\end{remark}

% ===========================================================================
\section{Volumetric mass-conservation theory}
\label{sec:theory}

\subsection{The physical mass invariant}
\label{sec:T1}

The conserved invariant of the incompressible/weakly-compressible flow is the
\emph{physical} mass
\begin{equation}
M \;=\; \int_{\Omega} \rho\,dV
     \;=\; \int_{\Omega_c} \rho\,|\mathcal{J}_{3}|\,d\xi\,d\eta\,d\zeta,
\label{eq:M-physical}
\end{equation}
where $\Omega\subset\mathbb{R}^3$ is the physical fluid domain and $\Omega_c$ its
logical pre-image, a uniform cuboid in $(\xi,\eta,\zeta)$.

A naive ``mass'' formed by summing the nodal density over computational indices
with \emph{unit} weight,
\begin{equation}
\widetilde{M}
\;=\;\sum_{p}\rho_p
\;\approx\;\int_{\Omega_c}\rho\,d\xi\,d\eta\,d\zeta
\;=\;\int_{\Omega}\frac{\rho}{|\mathcal{J}_{3}|}\,dV ,
\label{eq:Mtilde-wrong}
\end{equation}
is \emph{not} a discretization of \eqref{eq:M-physical}: it integrates over
computational space and therefore weights each physical sub-region by the inverse
local Jacobian $|\mathcal{J}_3|^{-1}$. On a stretched body-fitted mesh the metric
varies strongly---grid lines cluster near the wall and in the lee of the hill
where $|\mathcal{J}_3|$ is small---so $\widetilde{M}$ over-counts the
node-clustered (small-cell) regions and under-counts the coarse-cell regions.
Controlling $\widetilde{M}$ would drive the density field toward a spatially
\emph{biased} mean and inject a spurious, mesh-dependent body force into the
momentum equation. Only the genuinely physical functional \eqref{eq:M-physical},
weighted by $|\mathcal{J}_3|$, is the correct invariant to monitor and to
regulate.

\subsection{Quadrature of the physical mass and the mean density}
\label{sec:T2}

Discretize \eqref{eq:M-physical} by a nodal quadrature rule. Let $p=(i,j,k)$
index the unique interior lattice nodes and let $w_p>0$ be the \emph{physical
control volume} assigned to node $p$ (constructed in \S\ref{sec:weights}). The
mass and total volume are approximated by
\begin{align}
M &\;\approx\; \sum_{p} \rho_p\, w_p ,
\label{eq:M-quad}\\[2pt]
V &\;=\; \sum_{p} w_p
   \;\approx\; \int_{\Omega} dV \;=\; |\Omega| ,
\label{eq:V-quad}
\end{align}
and the \emph{volume-weighted mean density} is their ratio,
\begin{equation}
\langle\rho\rangle_V
\;\equiv\;\frac{M}{V}
\;=\;\frac{\sum_{p}\rho_p\,w_p}{\sum_{p}w_p}.
\label{eq:rho-mean}
\end{equation}
By construction $\langle\rho\rangle_V$ is the exact discrete average that equals
unity when $\rho_p\equiv 1$ for all $p$, \emph{independently of the mesh
stretching}, because the same weights $w_p$ appear in numerator and denominator.
This is the defining property the unit-weight surrogate \eqref{eq:Mtilde-wrong}
lacks. In the implementation the weighted sum $\sum_p \rho_p w_p$ is evaluated by
a GPU reduction over the unique interior node set $i\in[3,N_x{-}5]$,
$j\in[3,N_{y,d}{-}5]$, $k\in[3,N_z{-}4]$, summed across MPI ranks, and divided by
the precomputed global volume $V=\sum_p w_p$ to yield $\langle\rho\rangle_V$ on
the root rank. The fidelity of \eqref{eq:rho-mean} as a model of
\eqref{eq:M-physical} rests entirely on how faithfully $w_p$ reproduces the
physical control volume, constructed in \S\ref{sec:weights}.

\subsection{The correction operator: a discrete mass governor}
\label{sec:T4}

\subsubsection{Definition}
At every time step the solver forms $\langle\rho\rangle_V$ from
\eqref{eq:rho-mean} and defines the scalar correction
\begin{equation}
\delta\rho \;\equiv\; 1-\langle\rho\rangle_V ,
\label{eq:deltarho}
\end{equation}
the target reference density being unity in lattice units. The root rank computes
$\delta\rho$ and broadcasts it to all ranks. It is then applied as a
\emph{uniform} additive shift, injected entirely into the rest population $f_0$
at every interior node $p$:
\begin{equation}
\rho_p \;\leftarrow\; \rho_p + \delta\rho,
\qquad
f_{0,p} \;\leftarrow\; f_{0,p} + \delta\rho .
\label{eq:apply}
\end{equation}
We now establish three exact properties.

\subsubsection{(a) Mass change is exactly \texorpdfstring{$\delta\rho$}{drho} per node}
Since $\rho_p=\sum_q f_{q,p}$ and only $f_0$ is altered,
\begin{equation}
\Delta\rho_p \;=\; \sum_q \Delta f_{q,p}
\;=\; \Delta f_{0,p} \;=\; \delta\rho .
\label{eq:mass-change}
\end{equation}
The increment is identical at every node, so the change in the discrete
volume-weighted mass is
\begin{equation}
\Delta M
=\sum_p w_p\,\Delta\rho_p
=\delta\rho\sum_p w_p
=\delta\rho\,V .
\label{eq:DeltaM}
\end{equation}

\subsubsection{(b) Momentum is exactly invariant}
The D3Q19 rest velocity is $\mathbf{e}_0=(0,0,0)$. The momentum moment changes by
\begin{equation}
\Delta(\rho u_\alpha)_p
=\sum_q e_{q,\alpha}\,\Delta f_{q,p}
=e_{0,\alpha}\,\delta\rho
=0,
\qquad \alpha\in\{x,y,z\}.
\label{eq:momentum-invariant}
\end{equation}
Because the carrier population has zero lattice velocity, the injected mass
carries \emph{no} momentum: the correction lives in the kernel of every momentum
moment. This is precisely why the rest population $f_0$, and not any moving
population, is chosen as the carrier---any $\mathbf{e}_q\ne\mathbf 0$ would couple
the mass governor to the momentum equation and perturb the velocity field.

\subsubsection{(c) One-shot restoration of the unit mean}
Apply \eqref{eq:apply} and recompute the mean using \eqref{eq:rho-mean} and the
constancy of $\delta\rho$,
\begin{equation}
\langle\rho\rangle_V^{\,\mathrm{new}}
=\frac{\sum_p w_p(\rho_p+\delta\rho)}{\sum_p w_p}
=\frac{\sum_p w_p\rho_p}{\sum_p w_p}+\delta\rho
=\langle\rho\rangle_V+\bigl(1-\langle\rho\rangle_V\bigr)
=1 .
\label{eq:one-shot}
\end{equation}
A spatially uniform shift moves the volume-weighted mean by exactly its own
magnitude, so the target $\langle\rho\rangle_V=1$ is reached in a \emph{single}
application; no iteration is required. The only residual is the floating-point
rounding of the parallel reduction that forms $\langle\rho\rangle_V$.

\subsubsection{Control-theoretic interpretation}
Equations \eqref{eq:deltarho}--\eqref{eq:one-shot} constitute a discrete
proportional ``mass governor'' with \emph{unit feedback gain}. Writing
$\chi^{(n)}=\langle\rho\rangle_V$ at step $n$ before correction, the corrector
realizes the fixed-point map
\begin{equation}
\chi \;\longmapsto\; \chi + (1-\chi),
\label{eq:fixed-point}
\end{equation}
whose unique fixed point $\chi^{\star}=1$ is reached in one step with zero
contraction remainder (the error multiplier is $1-\mathrm{gain}=0$). The scheme
is unconditionally stable in this idealized sense: it neither under- nor
over-shoots, and introduces no oscillatory feedback mode. In the running solver
the correction is applied with a one-step lag---$\langle\rho\rangle_V$ is measured
at the end of step $n$ and the shift is applied at step $n{+}1$---so the governor
removes the compressibility drift accumulated over a single streaming/collision
update each step, holding the global mean density pinned at unity to
reduction-rounding precision throughout the simulation.

\subsection{Relation to the lattice-Boltzmann moments and forcing}
\label{sec:T5}

The corrector is consistent with the moment structure of the LBM. A perturbation
$\Delta f_q=\delta\rho\,\delta_{q0}$ has moment signature
\begin{equation}
\sum_q \Delta f_q=\delta\rho,
\qquad
\sum_q \mathbf{e}_q\,\Delta f_q=\mathbf{e}_0\,\delta\rho=\mathbf{0},
\label{eq:injection-moments}
\end{equation}
so the injection populates the density moment alone and lies in the common kernel
of all three momentum moments: it is a pure, momentum-silent density shift. This
is the moment-space statement of \eqref{eq:momentum-invariant}.

The mass governor is complementary to the body force. The solver applies Guo
forcing \cite{GuoZhengShi2002} with Hermite expansion order controlled at compile
time. At first order the forcing population is
\begin{equation}
F_q \;=\; W_q\,\bigl(3\,c_{q,y}\bigr)\,\mathcal{F},
\label{eq:guo1}
\end{equation}
(pure lattice-velocity weighting, with $\mathcal F$ the driving force and
$c_{q,y}=e_{q,y}$), whose density moment vanishes identically,
\begin{equation}
\sum_q F_q
=3\,\mathcal{F}\sum_q W_q\,c_{q,y}=0
\quad\Longleftrightarrow\quad
\sum_q W_q\,e_{q,y}=0 ,
\label{eq:guo1-massmoment}
\end{equation}
the identity holding by the lattice isotropy of the weights. Hence the
first-order body force injects \emph{no} net mass; the full second-order Guo
form, $F_q=W_q\,\mathcal F[\,3(c_{q,y}-u_y)+9(\mathbf c_q\!\cdot\mathbf u)c_{q,y}\,]$,
likewise carries zero density moment by the same weight identities. The
accumulated unit-mean drift that the governor removes is therefore \emph{not}
forcing-induced mass; it is the residual compressibility error of the
weakly-compressible LBM time stepping on the stretched curvilinear mesh. The
volume-weighted corrector \eqref{eq:deltarho}--\eqref{eq:apply} is the discrete
device that removes exactly this drift, in one momentum-preserving step per
iteration, against the physically correct invariant $M=\int_\Omega\rho\,dV$ of
\S\ref{sec:T1}.

% ===========================================================================
\section{Control-volume weights and the partition-of-unity identity}
\label{sec:weights}

\subsection{Cell volumes}
\label{sec:T3-cells}
The mesh cell $c=(i_c,j_c,k_c)$ is the hexahedron spanned by logical indices
$[i_c,i_c{+}1]\times[j_c,j_c{+}1]\times[k_c,k_c{+}1]$. By the block structure
\eqref{eq:jacobian-block}, its physical volume separates into the spanwise chord
times a cross-sectional area,
\begin{equation}
V_c \;=\; |\Delta x_{i_c}|\;A_c,
\qquad
\Delta x_{i_c}=x_{i_c+1}-x_{i_c}.
\label{eq:Vc-split}
\end{equation}
Two equivalent constructions of $A_c$ are provided; their detailed numerical
analysis is deferred to \S\ref{sec:quadrature}--\S\ref{sec:numerical-analysis}.

\paragraph{Method~0 (Shoelace / surveyor area).}
The four cross-section corners $(y_a,z_a)$, $a=0,\dots,3$, taken cyclically at
logical positions $(j_c,k_c),(j_c{+}1,k_c),(j_c{+}1,k_c{+}1),(j_c,k_c{+}1)$,
define a planar quadrilateral whose exact area is the Gauss/shoelace formula
\begin{equation}
A_c^{\mathrm{Shoe}}
\;=\;\tfrac12\Bigl|\sum_{a=0}^{3}\bigl(y_a z_{a+1}-z_a y_{a+1}\bigr)\Bigr|,
\qquad (a{+}1)\ \mathrm{mod}\ 4 .
\label{eq:shoelace}
\end{equation}
This is the \emph{exact} area of the straight-chord quadrilateral. When the $V_c$
of adjacent cells are summed, the line integrals along shared edges cancel
pairwise (opposite orientation), so the total telescopes to the exact area of the
piecewise-linear mesh polygon. Hence Method~0 reproduces the discrete
(piecewise-linear) mesh volume \emph{exactly}; its only error is the $O(h^2)$
geometric chording of the truly curved boundary.

\paragraph{Method~1 (Jacobian, $3\times3$ Gauss--Legendre).}
Alternatively the area is the metric integral
\begin{equation}
A_c^{\mathrm{Jac}}
=\!\int_{j_c}^{j_c+1}\!\!\int_{k_c}^{k_c+1}\!\! J_{2D}(\xi,\zeta)\,d\xi\,d\zeta
\;\approx\;\sum_{a=1}^{3}\sum_{b=1}^{3} W_a W_b\,
J_{2D}\!\bigl(j_c+\xi_a,\,k_c+\zeta_b\bigr),
\label{eq:Ajac}
\end{equation}
a tensor-product $3$-point Gauss--Legendre (GL-3) rule. On the reference cell
$[0,1]$ the nodes and weights are
\begin{equation}
\{\xi_a\}=\Bigl\{\tfrac{1-\sqrt{3/5}}{2},\ \tfrac12,\ \tfrac{1+\sqrt{3/5}}{2}\Bigr\},
\qquad
\{W_a\}=\Bigl\{\tfrac{5}{18},\ \tfrac{8}{18},\ \tfrac{5}{18}\Bigr\}.
\label{eq:GL3}
\end{equation}
GL-$n$ integrates polynomials of degree $\le 2n{-}1$ exactly; for $n=3$ the rule
is exact through degree $5$ per direction \cite{StoerBulirsch2002,DavisRabinowitz1984}.
Because the nodes \eqref{eq:GL3} are off-grid, $J_{2D}$ at a Gauss point is
reconstructed by a $6$-point (degree-$5$) tensor Lagrange interpolation of nodal
$J_{2D}$ values,
\begin{equation}
J_{2D}(\xi,\zeta)\;\approx\;\sum_{m=0}^{5}\sum_{n=0}^{5}
\ell_m(\xi)\,\ell_n(\zeta)\,J_{2D}\bigl(s_j{+}m,\,s_k{+}n\bigr),
\qquad
\ell_m(\xi)=\!\!\prod_{\substack{r=0\\ r\ne m}}^{5}\frac{\xi-(s+r)}{(s{+}m)-(s{+}r)},
\label{eq:lagrange6}
\end{equation}
with a centered, wall-clamped stencil start $s$. The reconstruction is $O(h^6)$
in $J_{2D}$, exactly matching the degree-$5$ exactness of GL-3, so the
interpolation does \emph{not} lower the quadrature order. If the $6$-wide stencil
is unavailable (domain edge) or the reconstructed $J_{2D}$ is non-finite or
$\le 0$, the cell falls back to the Shoelace value \eqref{eq:shoelace}. Either way
the nodal cell volume is $V_c=|\Delta x_{i_c}|A_c$, and a verification routine
(\texttt{VerifyPhysicalDomainVolume3D}) checks the assembled weight-sum against the
discrete mesh volume: Method~0 reproduces it to relative error below $10^{-12}$
(printing \texttt{PASS}), whereas Method~1 \emph{deliberately} departs by the
$O(h^2)$ boundary-curvature term and the routine prints a \texttt{WARNING}. The
routine only prints \texttt{PASS}/\texttt{WARNING}; it does not assert or abort.

\subsection{Median-dual nodal weight}
\label{sec:T3-weight}
The control volume of node $p$ is the median-dual (vertex-centered finite-volume)
cell \cite{Barth1991}: each of the (at most eight) mesh cells touching $p$
contributes one eighth of its volume,
\begin{equation}
w_p \;=\; \frac{1}{8}\sum_{c\,\ni\,p} V_c ,
\label{eq:median-dual}
\end{equation}
the factor $1/8$ being the fraction of a hexahedral cell nearest one of its eight
corners. Two boundary adaptations are made without altering the weight's meaning:
(i) in the periodic spanwise direction the wrap is closed by mapping the first
interior layer's left cell $i=3$ to the cell index $N_x{-}5$, so the two physical
cells adjacent to the periodic seam are both counted exactly once; (ii) at a
solid wall in $k$ a node has cells on one side only, so the inner sum runs over
the available (one) cell rather than two.

\subsection{Partition-of-unity proof}
\label{sec:T3-pou}
\begin{proposition}[Exact volume closure]
\label{prop:pou}
With $w_p$ defined by \eqref{eq:median-dual}, the nodal weights sum to the total
mesh volume,
\begin{equation}
\sum_{p} w_p \;=\; \sum_{c} V_c \;=\; V_{\mathrm{mesh}} .
\label{eq:pou}
\end{equation}
\end{proposition}
\begin{proof}
Substitute \eqref{eq:median-dual} and exchange the order of summation,
\begin{equation}
\sum_{p} w_p
=\frac{1}{8}\sum_{p}\ \sum_{c\,\ni\,p} V_c
=\frac{1}{8}\sum_{c} V_c\,\bigl|\{\,p : p\ \text{is a corner of}\ c\,\}\bigr| .
\label{eq:pou-swap}
\end{equation}
Every interior hexahedral cell $c$ has exactly eight corner nodes, so the
incidence count is $|\{p:\,p\in c\}|=8$ and the prefactor $1/8$ cancels it
identically, giving $\sum_p w_p=\sum_c V_c$. The pairing of the $\tfrac18$ shares
with the $8$ corners is a telescoping over the cell--node incidence relation and
is therefore exact in floating point up to the order of summation; it does not
depend on cell shape, stretching, or curvature. This closure is an
exact-arithmetic identity for a periodic mesh whose seam is \emph{bit-closed};
the production streamwise grid seam is in fact open by $\sim\!3\times10^{-6}$,
which adds a small $\sim\!10^{-7}$ geometric residual to \eqref{eq:pou},
measured and decomposed in \S\ref{sec:E4}.

The boundary adaptations preserve the count. For the periodic spanwise seam,
remapping $i=3\to N_x{-}5$ relabels which cell supplies a corner share but leaves
each physical cell with exactly its eight one-eighth contributions, so no cell is
double-counted or dropped and the incidence count stays $8$. At a solid wall, a
boundary node legitimately borders fewer cells; the missing neighbours lie
outside $\Omega$ and contribute no volume, so the telescoping still assigns each
\emph{existing} cell's eighth-shares to its existing corner nodes. In every case
each cell contributes its full $V_c$ exactly once, proving \eqref{eq:pou}.
\end{proof}
Consequently the denominator in \eqref{eq:rho-mean}, $V=\sum_p w_p$ (assembled by
an MPI all-reduce into the global control volume), equals the discrete mesh
volume in exact arithmetic (to rounding, and exactly so only on a bit-closed
grid; see the refinement in \S\ref{sec:self-consistency}), and
$\langle\rho\rangle_V$ is a true volume average.

% ===========================================================================
\section{Cell-volume quadrature: Shoelace vs.\ Jacobian GL-3}
\label{sec:quadrature}

Both cell-volume estimators are quadratures of the \emph{same} area integral
\begin{equation}
  A(j,k) \;=\; \int_{j}^{\,j+1}\!\!\int_{k}^{\,k+1}
  J_{2D}(\xi,\zeta)\,\mathrm{d}\xi\,\mathrm{d}\zeta ,
  \qquad
  V_{\mathrm{cell}}(i,j,k)=\lvert\Delta x_i\rvert\,A(j,k),
  \label{eq:na-Adef}
\end{equation}
and differ only in the order to which they resolve the curvature of $J_{2D}$.

\subsection{N1: two quadratures of the same area integral}
\label{sec:na-N1}

\paragraph{Method~0 --- Shoelace (exact straight-chord area).}
The first estimator (\texttt{MassCorrectionCellVolume}) replaces the curved cell
boundary by its four straight chords and computes the exact area of the resulting
planar quadrilateral with the surveyor (Gauss/shoelace) formula \eqref{eq:shoelace}.
Equation~\eqref{eq:shoelace} is the exact integral of $J_{2D}$ for a \emph{bilinear}
map whose edges are straight, i.e.\ for the piecewise-linear (``discrete'') mesh.
When adjacent quadrilaterals are summed, every interior edge contributes with
opposite sign in the two cells that share it and cancels; the sum therefore
telescopes to the exact area of the discrete-mesh polygon. As a consequence
Method~0 reproduces the discrete-mesh volume to machine precision \emph{on a
bit-closed periodic mesh} (this is precisely what
\texttt{VerifyPhysicalDomainVolume3D} checks; see \S\ref{sec:na-N4}); on the
production grid the streamwise periodic seam is not bit-closed and contributes a
$\sim\!10^{-7}$ geometric residual that exceeds the verifier's $10^{-12}$ pass
gate, quantified in \S\ref{sec:E4}. However, the chords are an $O(h^2)$ approximation of the true
curved boundary, so Method~0 carries a fixed \emph{geometric} discretisation error
\begin{equation}
  A^{\mathrm{Shoe}}(j,k) - A(j,k) = O(h^2)\cdot \kappa ,
  \label{eq:na-shoe-error}
\end{equation}
where $\kappa$ measures the local boundary curvature. This error does not vanish
as the quadrature is refined per cell --- it is a property of the chorded domain,
not of the integration rule.

\paragraph{Method~1 --- Jacobian $3\times3$ Gauss--Legendre.}
The second estimator (\texttt{ComputeJacobianMassCorrectionWeights}) integrates
the in-plane metric $J_{2D}$ with the tensor-product GL-3 rule
\eqref{eq:Ajac}--\eqref{eq:GL3}. The decisive point for the error analysis is that
the solver does \emph{not} feed a closed-form $J_{2D}$ to the quadrature: the Gauss
abscissae are off-grid, so at every Gauss point $J_{2D}$ is first \emph{reconstructed}
from nodal samples by the $6$-point (degree-$5$) tensor Lagrange interpolation of
\S\ref{sec:na-N2}, and GL-3 then integrates \emph{that} reconstruction. A 1D
$n$-point Gauss rule is exact for polynomials of degree $\le 2n-1$; with $n=3$ this is
degree~5 per direction, and the reconstruction is itself a degree-5 polynomial per
direction. \emph{GL-3 therefore integrates the degree-5 reconstruction exactly}: the
Gauss remainder is identically zero on the quantity actually integrated, so the
quadrature contributes \emph{no} error to leading order. The convergence-limiting
error of Method~1 is consequently the degree-5 \emph{Lagrange reconstruction} error of
the sampled metric,
\begin{equation}
  A^{\mathrm{Jac}}(j,k)-A(j,k)
  \;=\; \int_{j}^{j+1}\!\!\int_{k}^{k+1}\!\!
        \bigl(\tilde J_{2D}-J_{2D}\bigr)\,\mathrm{d}\xi\,\mathrm{d}\zeta
  \;=\; O(h^{6})\,J_{2D}^{(6)},
  \label{eq:na-method1-error}
\end{equation}
i.e.\ the per-cell interpolation defect \eqref{eq:na-interp-error} integrated over the
cell. Summed over the $O(h^{-2})$ cells of the cross-section the global area error of
Method~1 is $O(h^{6})$---interpolation-limited, not quadrature-limited---versus the
$O(h^{2})$ geometric chord bias of Method~0. This sixth-order rate is confirmed
empirically in Section~\ref{sec:experiments} (the measured orders are $\approx 6.0$
down to the round-off floor).

\paragraph{Scope of the Gauss remainder.}
For completeness we record the classical 1D Gauss remainder
\begin{equation}
  R_n \;=\; \frac{(b-a)^{2n+1}\,(n!)^4}
                 {(2n+1)\,\bigl((2n)!\bigr)^3}\,
            f^{(2n)}(\theta),\qquad \theta\in(a,b),
  \label{eq:na-gauss-remainder}
\end{equation}
which for $n=3$ on the computational cell $b-a=1$ gives the closed-form constant
\begin{equation}
  R_3 = \frac{(3!)^4}{7\,\bigl(6!\bigr)^3}\,J_{2D}^{(6)}(\theta)
      = \frac{1296}{7\cdot 373{,}248{,}000}\,J_{2D}^{(6)}(\theta)
      = \frac{1296}{2{,}612{,}736{,}000}\,J_{2D}^{(6)}(\theta)
      \approx 4.96\times10^{-7}\,J_{2D}^{(6)}(\theta).
  \label{eq:na-R3}
\end{equation}
Equation~\eqref{eq:na-R3} bounds the quadrature-\emph{only} error that would arise if
a general smooth $J_{2D}$ with degree-$>5$ content were fed \emph{directly} to GL-3.
In the solver it is \emph{not} the operative error: the integrand is the degree-5
reconstruction, on which $R_3\equiv0$. It is retained only as the bound on the (here
null) quadrature contribution; the binding term is the interpolation error
\eqref{eq:na-method1-error}. The implementation falls back to \eqref{eq:shoelace} for
any cell where the 6-wide interpolation stencil is unavailable or where the
reconstructed $J_{2D}$ is non-finite or $\le0$, guaranteeing a strictly positive,
finite weight everywhere.

\subsection{N2: degree-5 Lagrange reconstruction at the Gauss points}
\label{sec:na-N2}

The Gauss abscissae $j+\xi_a$, $k+\zeta_b$ are non-nodal, so $J_{2D}$ must be
reconstructed there from the nodal metric. The code
(\texttt{InterpolateJ2D\_Lagrange6}) uses the tensor product
\eqref{eq:lagrange6}. The interpolation error of a degree-5 reconstruction is
\begin{equation}
  J_{2D}-\tilde{J}_{2D}
  = \frac{J_{2D}^{(6)}(\theta)}{6!}\prod_{r=0}^{5}\bigl(\xi-(s+r)\bigr)
  = O(h^6)\,J_{2D}^{(6)} .
  \label{eq:na-interp-error}
\end{equation}
This interpolation defect is the \emph{binding} (dominant) error term of Method~1.
GL-3 has degree of exactness~5 and the reconstruction \eqref{eq:lagrange6} is a
degree-5 polynomial, so the quadrature \eqref{eq:Ajac} integrates the reconstruction
\emph{exactly} and adds nothing to leading order---the Gauss remainder
\eqref{eq:na-R3} is null on the reconstruction. What survives is precisely the gap
between the reconstruction and the true metric, $\tilde{J}_{2D}-J_{2D}=O(h^6)$ of
\eqref{eq:na-interp-error}; integrated over each cell and summed it sets the global
$O(h^6)$ rate \eqref{eq:na-method1-error}. Method~1 thus integrates a degree-5
polynomial reconstruction of the smooth metric, capturing the sub-cell curvature that
the straight chord of Method~0 ignores, and its accuracy is limited not by the
quadrature but by how well that degree-5 reconstruction approximates $J_{2D}$, i.e.\
by $J_{2D}^{(6)}h^6$.

\paragraph{Stencil placement.}
\texttt{SelectStencilStart} centres the 6-point window on the target cell by
choosing $s=\mathrm{cell}-2$ and clamping to $[\,\text{lo},\,\text{hi}-5\,]$, so
that interior cells receive a symmetric stencil while near-boundary cells degrade
gracefully to a one-sided stencil rather than reaching past the valid data. In the
wall-normal direction the $\zeta$-stencils are restricted to physical nodes
($k=3,\dots,N_z{-}4$): the buffer rows $k=2$ and $k=N_z{-}3$ hold metric values
that are not part of the physical domain, so near-wall GL cells use the one-sided
physical stencils $k=3..8$ and $k=N_z{-}9..N_z{-}4$. Because an interior cell's
centred stencil reaches two rows into the neighbouring MPI subdomain in the
streamwise ($\xi$) direction, the ghost rows of $J_{2D}$ ($j=0,1,2$ and
$j=N_y{-}3,\dots,N_y{-}1$) must be exchanged across ranks before the quadrature
loop. The solver performs exactly this halo exchange of $J_{2D}$ prior to invoking
Method~1; without it, partition-boundary cells would silently fall back to the
lower-order shoelace estimate.

% ===========================================================================
\section{Numerical analysis and error budget}
\label{sec:numerical-analysis}

\subsection{N3: floating-point exactness of the correction}
\label{sec:na-N3}

\paragraph{Nodal control volumes and partition of unity.}
Each interior node $p$ is assigned the median-dual weight
$w_p=\tfrac18\sum_{c\ni p}V_{\mathrm{cell}}(c)$ of \eqref{eq:median-dual}, with the
periodic-$x$ wrap and one-sided near-wall stencils handled explicitly. Since each
cell is shared by its eight corner nodes, the weights form a partition of unity
that is exact in floating point up to summation order
(Proposition~\ref{prop:pou}),
\begin{equation}
  \sum_p w_p = \sum_c V_{\mathrm{cell}}(c) = V ,
  \label{eq:na-pou}
\end{equation}
where $V$ is obtained once at start-up by \texttt{MPI\_Allreduce}.

\paragraph{One-shot algebraic exactness.}
At every step the solver forms $\langle\rho\rangle_V=\tfrac1V\sum_p w_p\rho_p$ and
$\delta\rho=1-\langle\rho\rangle_V$, then applies the uniform shift
$\rho\leftarrow\rho+\delta\rho$ together with $f_0\leftarrow f_0+\delta\rho$ at
every interior node. Because the lattice rest velocity is
$\mathbf{e}_0=\mathbf{0}$, the shift injects mass but \emph{no} momentum
(\eqref{eq:momentum-invariant}). Applying the uniform shift and using
\eqref{eq:na-pou} gives
\begin{equation}
  \langle\rho\rangle_V^{\text{new}}
  = \frac{1}{V}\sum_p w_p\,(\rho_p+\delta\rho)
  = \langle\rho\rangle_V + \delta\rho\,\frac{\sum_p w_p}{V}
  = \langle\rho\rangle_V + \delta\rho = 1 ,
  \label{eq:na-oneshot}
\end{equation}
so the mean density is restored to unity in a \emph{single} application; the only
residual is the floating-point rounding of the reduction.

\paragraph{Conditioning of the reduction.}
The reduced quantity $S=\sum_p w_p\rho_p$ is a sum of strictly positive terms
($w_p>0$, $\rho_p>0$), so there is no catastrophic cancellation: the condition
number $\bigl(\sum_p\lvert w_p\rho_p\rvert\bigr)/\lvert\sum_p w_p\rho_p\rvert=1$.
A naive sequential accumulation has a worst-case relative error $O(n\varepsilon)$.
The device reduction is, however, only \emph{two-level}: the kernel
(\texttt{ReduceRhoWeightedSum\_Kernel}) performs a balanced binary-tree reduction
\emph{within} each $B{=}256$-thread block, after which the $\sim n/B$ block partials
are summed \emph{sequentially} on the host
(\texttt{ComputeVolumeWeightedRhoAverageRoot}:
\texttt{for(b) rho\_LocalWeightedMass += rho\_partial\_h[b]}), followed by a
cross-rank \texttt{MPI\_Reduce}. The intra-block tree therefore lowers the rounding
floor only to $O(\varepsilon\,n/B)$ with $B{=}256$---\emph{not} to the full
$O(\varepsilon\log_2 n)$ of a single balanced tree. For $n\sim10^6$ nodes in double
precision,
\begin{equation}
  n\varepsilon \approx 2.2\times10^{-10},
  \qquad
  \frac{n}{B}\,\varepsilon \approx \frac{10^{6}}{256}\,\varepsilon \approx 9\times10^{-13},
  \label{eq:na-rounding}
\end{equation}
i.e.\ the intra-block tree buys roughly a factor $B=256$ over the naive sequential
sum (a realistic floor $\sim\!10^{-12}$ at production scale), well short of the
$\sim\!4\times10^{-15}$ that a single fully balanced tree would reach. The remaining
caveat is determinism: floating-point addition is non-associative, so the
across-rank \texttt{MPI\_Reduce} is bit-reproducible only for a fixed rank order
and fixed local reduction tree. The implementation preserves both, so a given run
is reproducible, while a change in rank count perturbs $\delta\rho$ at the
$O(\varepsilon\,n/B)$ level seen in \eqref{eq:na-rounding}.

\paragraph{Mass governor interpretation.}
With \texttt{FORCE\_HERMITE\_ORDER}${=}2$ (full Guo forcing) or its first-order
variant, the body force carries zero density moment (\eqref{eq:guo1-massmoment}),
so it injects no net mass; the residual compressibility drift accumulated by the
LBM time stepping is what the corrector removes. Equation \eqref{eq:na-oneshot} is
therefore a discrete mass governor with unit feedback gain and a one-step lag,
whose closed-loop fixed point is $\langle\rho\rangle_V=1$.

\subsection{N4: why the curved domain prefers Method~1}
\label{sec:na-N4}

It is tempting to call Method~0 ``exact'' and stop there, but the two estimators
are exact for \emph{different} domains. Method~0 is exact for the chorded,
piecewise-linear mesh; Method~1 converges to the true curved physical domain as
the metric is resolved. For a periodic hill whose lower wall is a smooth curve,
the chorded domain is \emph{not} the physical domain, and the $O(h^2)$ bias of
\eqref{eq:na-shoe-error} contaminates every control volume $w_p$ and hence the
weighted mean \eqref{eq:rho-mean}. Method~1, by integrating the smooth mapping
\eqref{eq:jacobian-block}, produces weights that converge to the true physical
control volumes; this is why the default is \texttt{CELL\_VOLUME\_METHOD}${=}1$.

This also explains the design of \texttt{VerifyPhysicalDomainVolume3D}. The check
compares the summed weights against $V_{\mathrm{discrete\,mesh}}$, the exact
discrete-mesh volume formed from unique shoelace cells. For Method~0 the relative
difference is required to pass the $<10^{-12}$ tolerance --- it \emph{is} the
discrete-mesh volume (telescoping, \S\ref{sec:na-N1}). For Method~1 the routine is
expected to report a \emph{nonzero} difference of the order of the
boundary-curvature term \eqref{eq:na-shoe-error}: Method~1 deliberately departs
from the chorded volume to recover the curved physical volume, and the routine
additionally prints the analytic
$L_X\bigl(L_Y L_Z - \int_0^{L_Y} h(y)\,\mathrm{d}y\bigr)$ reference (Simpson,
$N{=}10^4$) so that the discrete-vs-analytic gap can be read off directly. A
``WARNING: differs from discrete mesh volume'' for Method~1 is therefore the
intended, physically correct behaviour, not a failure.

\subsection{Error budget}
\label{sec:na-budget}

Table~\ref{tab:na-budget} collects the leading error contributions. The geometric
bias of Method~0 dominates and does not shrink under per-cell quadrature refinement
(it converges only under mesh refinement, at the $O(h^2)$ rate); Method~1 removes it
down to the degree-5 Lagrange reconstruction error $O(h^6)$ of
\eqref{eq:na-method1-error} (the binding term---GL-3 integrates that reconstruction
exactly, so the Gauss remainder \eqref{eq:na-R3} is non-binding here), which in turn
falls to the floating-point floor of the parallel reduction by the finest grids. The
corrector itself is algebraically exact \eqref{eq:na-oneshot}, so the overall
mass-closure error is set by the conditioning analysis of \S\ref{sec:na-N3}.

\begin{table}[t]
  \centering
  \caption{Error budget for the curvilinear mass-conservation corrector.
  $h$ is the cell size, $n\sim10^6$ the number of interior reduction nodes,
  $\varepsilon=2.2\times10^{-16}$ the double-precision unit roundoff. ``Per cell''
  errors refer to the area integral $A(j,k)$ of \eqref{eq:na-Adef}; ``Global''
  aggregates over the $O(h^{-2})$ cross-section cells or the $n$ nodes.}
  \label{tab:na-budget}
  \small
  \resizebox{\textwidth}{!}{%
  \begin{tabular}{@{}llll@{}}
    \toprule
    Source & Mechanism & Per-cell order & Global magnitude \\
    \midrule
    Method~0 chord vs.\ curve     & geometric, $O(h^2)$ boundary chord & $O(h^2)\kappa$ & $O(h^2)$ chord bias (mesh-converging) \\
    Lagrange-6 reconstruction (\textbf{binding}) & degree-5 interp.\ \eqref{eq:na-interp-error} & $O(h^6)\,J^{(6)}_{2D}$ & $O(h^6)$ --- \textbf{dominant}, sets the global rate \\
    Method~1 GL-3 quadrature      & Gauss remainder \eqref{eq:na-R3}   & exact on the deg-5 reconstr. & non-binding (null on the reconstruction) \\
    Partition-of-unity closure    & $\sum_p w_p = V$ \eqref{eq:na-pou} & exact (telescoping) & $O(\varepsilon\,n/B)$ \\
    Correction algebra            & one-shot shift \eqref{eq:na-oneshot} & exact & exact ($\to\langle\rho\rangle_V=1$) \\
    Reduction (naive sequential)  & accumulation roundoff               & --- & $O(n\varepsilon)\approx2.2\times10^{-10}$ \\
    Reduction (two-level: block tree + host seq.) & intra-block tree, then sequential host sum & --- & $O(\varepsilon\,n/B)\approx9\times10^{-13}$ ($B{=}256$) \\
    Cross-rank \texttt{MPI\_Reduce} & non-associative FP, fixed order   & --- & reproducible; $O(\varepsilon\,n/B)$ jitter on rank-count change \\
    \bottomrule
  \end{tabular}}
\end{table}

\noindent
In summary, Method~1 replaces a fixed $O(h^2)$ geometric modelling error by an
interpolation-limited $O(h^6)$ error (the degree-5 Lagrange reconstruction, which
GL-3 integrates exactly) several orders of magnitude smaller and convergent under
mesh refinement, while the median-dual partition of unity and the momentum-preserving
rest-population shift make the resulting mass corrector algebraically exact, with a
residual governed only by the well-conditioned, tree-summed parallel reduction.

% ===========================================================================
\section{Numerical experiments}
\label{sec:experiments}

The theory of \S\ref{sec:theory}--\S\ref{sec:numerical-analysis} is validated by
an \emph{independent, host-only verification harness} (a standalone C++ program,
compiled with g++ 11.5.0 and \texttt{-O2 -std=c++17}, that re-implements the
solver-identical quadrature, interpolation, and corrector arithmetic). The harness
exercises three checks: (E1) polynomial exactness of the GL-3 rule, (E2)
$h$-refinement convergence of the two area estimators on a curved, genuinely
two-dimensional (non-separable) curvilinear map, and (E3) one-shot mass
restoration with momentum preservation. It
is deliberately decoupled from the GPU solver so that the mathematical properties
are tested in isolation from the production code path.

\subsection{E1: GL-3 polynomial exactness}

\begin{table}[t]
\centering
\caption{Polynomial exactness of the tensor 3-point Gauss--Legendre rule
($\int_0^1\xi^m\,d\xi$) used in the cell-volume scheme. Machine-zero error for
$m\le5$ confirms the $2n-1=5$ degree of exactness; the $m=6$ error matches the
analytic Gauss remainder $C\,f^{(6)}$ with $C=(3!)^4/[7\,(6!)^3]=4.96\times10^{-7}$.}
\label{tab:gl3_exactness}
\begin{tabular}{c r r r}
\toprule
$m$ & {Exact $1/(m{+}1)$} & {GL-3 quadrature} & {Abs.\ error} \\
\midrule
0 & 1.0000000000000000 & 1.0000000000000000 & 0.000e+00 \\
1 & 0.5000000000000000 & 0.5000000000000000 & 0.000e+00 \\
2 & 0.3333333333333333 & 0.3333333333333334 & 5.551e-17 \\
3 & 0.2500000000000000 & 0.2500000000000000 & 0.000e+00 \\
4 & 0.2000000000000000 & 0.2000000000000000 & 0.000e+00 \\
5 & 0.1666666666666667 & 0.1666666666666667 & 2.776e-17 \\
6 & 0.1428571428571428 & 0.1425000000000000 & 3.571e-04 \\
7 & 0.1250000000000000 & 0.1237500000000000 & 1.250e-03 \\
\bottomrule
\end{tabular}
\end{table}

Table~\ref{tab:gl3_exactness} confirms the tensor GL-3 rule is exact to machine
precision for polynomials of degree $\le5$ and first errs at degree~6 by
$3.571\times10^{-4}$, which reproduces the analytic Gauss remainder $C\,f^{(6)}$
(with $f=\xi^6$, $f^{(6)}=720$, $C\cdot720=3.571\times10^{-4}$) to all printed
digits. The quadrature therefore has exactly the claimed $2n-1=5$ degree of
exactness.

\subsection{E2: \texorpdfstring{$h$}{h}-refinement convergence}

\begin{table}[t]
\centering
\caption{Grid-convergence of the cross-section area on a genuinely two-dimensional,
\emph{non-separable} curvilinear $y$--$z$ map: $y=\xi$,
$z=\zeta\,(1-a\sin\pi\xi)+(C-D\sin\pi\xi)(1-\cos\pi\zeta)/\pi$, giving the in-plane
Jacobian $J_{2D}=1-a\sin(\pi\xi)+C\sin(\pi\zeta)-D\sin(\pi\xi)\sin(\pi\zeta)$
($a=0.30,\,C=0.20,\,D=0.10$; $J_{2D}\in[0.7,1.2]>0$ everywhere; exact area
$1-0.2/\pi-0.4/\pi^2=0.8958095493$). The $\xi\zeta$ cross term makes $J_{2D}$ vary in
\emph{both} directions, so the wall-normal Lagrange interpolation is genuinely
exercised. Method~0 is the straight-chord Shoelace area; Method~1 is the Jacobian
$3\times3$ Gauss--Legendre quadrature of the $6$-point Lagrange interpolation of
\emph{sampled} nodal $J_{2D}$ (solver-identical). Observed order
$p=\log_2(e_{h}/e_{h/2})$.}
\label{tab:gl3_convergence}
\begin{tabular}{c c r c r c}
\toprule
$N$ & $h$ & {Method 0 error} & {order} & {Method 1 error} & {order} \\
\midrule
  8 & 0.125000 & 2.983e-03 & {--}   & 2.636e-06 & {--}   \\
 16 & 0.062500 & 7.443e-04 & 2.00   & 3.053e-08 & 6.43   \\
 32 & 0.031250 & 1.860e-04 & 2.00   & 4.268e-10 & 6.16   \\
 64 & 0.015625 & 4.649e-05 & 2.00   & 6.464e-12 & 6.04   \\
128 & 0.007813 & 1.162e-05 & 2.00   & 1.040e-13 & 5.96   \\
256 & 0.003906 & 2.905e-06 & 2.00   & 1.887e-15 & {$\dagger$} \\
\bottomrule
\end{tabular}
\par\smallskip
\footnotesize $\dagger$ At $N{=}256$ Method~1 has reached the
$\sim\!2\times10^{-15}$ floating-point (reduction-rounding) floor; the error has
dropped into the round-off band, so the apparent ``order'' there is no longer a
genuine convergence rate. $N{=}128$ ($1.040\times10^{-13}$, order $5.96$) is still a
genuine sixth-order step, so saturation first appears at $N{=}256$.
\end{table}

Table~\ref{tab:gl3_convergence} cleanly separates the two schemes. The test map is
\emph{non-separable}---its in-plane Jacobian $J_{2D}$ carries a genuine $\xi\zeta$
cross term and varies in both the streamwise ($\xi$) and wall-normal ($\zeta$)
directions---so the tensor $6$-point Lagrange reconstruction is genuinely exercised
in the wall-normal direction as well, not merely streamwise (a separable $J_{2D}(\xi)$
would make the $\zeta$-interpolation trivially exact and leave that path untested).
The straight-chord Shoelace (Method~0) converges at a textbook $O(h^2)$ (observed
order $2.00$ at every level) because its piecewise-linear facets chord the curved
(one-signed) top wall (an Euler--Maclaurin trapezoid error), whereas the Jacobian
GL-3 scheme (Method~1) converges at the predicted sixth order (observed orders
$6.43,6.16,6.04,5.96$ from $N{=}8$ to $N{=}128$). The binding error is the degree-5
\emph{Lagrange reconstruction} of the sampled metric, which GL-3 integrates exactly,
so the reconstruction---not the quadrature---sets the $O(h^6)$ rate. At $N{=}256$ the
Method~1 error has fallen to $1.887\times10^{-15}$, i.e.\ into the
$\sim\!2\times10^{-15}$ reduction-rounding floor (marked $\dagger$); $N{=}128$
($1.040\times10^{-13}$, order $5.96$) is still a genuine sixth-order step, so the
saturation first appears at $N{=}256$, one level beyond the last clean rate. Because
the analytic map keeps $J_{2D}\in[0.7,1.2]$ strictly positive everywhere, the
solver's \texttt{isfinite}/${\le}0$ Shoelace fallback is never triggered in this test
(fallback count $=0$): it is a never-exercised safety path here, which we report
honestly as untested rather than claim as validated. This confirms the error-budget
claim that Method~1 replaces a fixed $O(h^2)$ geometric bias with a convergent,
interpolation-limited $O(h^6)$ error several orders of magnitude smaller.

\subsection{E3: one-shot mass restoration and momentum preservation}

On a $16\times16$ curvilinear patch, median-dual weights were built with each
volume method, the nodal density was perturbed, and the correction
$\delta\rho=1-\langle\rho\rangle_V$ was injected into the rest population $f_0$.
For both methods the perturbed mean $\langle\rho\rangle_V=0.99700000$
($\delta\rho=+3.0\times10^{-3}$) returned to unity after a single application to
within the reduction-rounding floor: Method~0 (Shoelace, $V=0.8965538$) gave
$|1-\langle\rho\rangle_V^{\mathrm{new}}|=1.110\times10^{-15}$ and Method~1 (Jacobian
GL-3, $V=0.8958096$, agreeing with the analytic area $0.8958095$ to the
$\sim\!3\times10^{-8}$ Method-1 area error at this $N{=}16$ resolution) gave
$2.220\times10^{-16}$. The Method-1 path again recorded zero Shoelace fallbacks (the
smooth $J_{2D}>0$ never trips the \texttt{isfinite}/${\le}0$ guard). A separate
momentum test injected $\delta\rho=0.0123456789$ into $f_0$ of a D3Q19 stencil
($\mathbf e_0=\mathbf0$): the density rose by exactly $\delta\rho$ while the momentum
$\sum_q\mathbf e_q f_q$ stayed at $(0,0,0)$ \emph{exactly}. This validates the
one-shot algebraic exactness \eqref{eq:na-oneshot} and confirms that the rest
population is the correct, momentum-silent mass carrier. The harness artefacts reside
under \path{docs/numexp/} (\path{verify_gl3.cpp}, \path{gl3_results.txt}).

% ---------------------------------------------------------------------------
\subsection{E4: partition-of-unity on the production grid --- a bit-exactness audit}
\label{sec:E4}

Experiments E1--E3 exercise the quadrature on analytic patches. E4 closes the
loop on the \emph{actual} $\mathrm{Re}_h{=}2800$ production grid
(\path{J_Frohlich/adaptive_3.fine grid_I513_J257_s0.950000.dat};
$N_x{\times}N_y{\times}N_z=257{\times}513{\times}257$, hence
$256{\times}512{\times}256=3.36\times10^{7}$ cells) and answers one question with
hard data: is the start-up control volume $\Sigma_w=\sum_p w_p$ (the
$\tfrac18$-split dual weights of \texttt{InitializeMassCorrectionWeights})
\emph{bit-exactly} equal to the discrete main-field volume
$V_{\mathrm{disc}}=\sum_c V_c$ (\texttt{ComputeGlobalDiscreteShoelaceVolume3D})?
A self-contained host harness (\path{docs/numexp/test_volume_partition.cpp},
single process, $jp{=}1$ whole-domain so the per-rank slice equals the global
array) reads the real grid through a byte-faithful copy of the solver's
\texttt{GenerateMesh\_X}\,/\,\texttt{ReadExternalGrid\_YZ} and evaluates both sums
in the solver's exact loop order and with its exact constants ($0.125$, the
Shoelace formula). The full transcript is in
\path{docs/numexp/partition_results.txt}.

\paragraph{The four bit-exactness metrics (Method~0).}
On the production grid,
\begin{equation}
\begin{aligned}
V_{\mathrm{disc}} &= 1.14358939042152798\times 10^{2}, &
\Sigma_w &= 1.14358953365302284\times 10^{2},\\[2pt]
\text{(i) exact }(==) &: \ \textbf{FALSE}, &
\text{(ii) ULP gap} &= 1{,}007{,}902{,}042 \approx 1.0\times10^{9},\\[2pt]
\text{(iii) }|\Sigma_w-V_{\mathrm{disc}}| &= 1.432\times10^{-5}, &
\text{(iv) rel} &= 1.252\times10^{-7}.
\end{aligned}
\label{eq:E4-metrics}
\end{equation}
The two doubles are \textbf{not} bit-identical (\textsc{bit-exact: no}), and,
decisively, the gap ($\sim\!10^{-7}$ relative, $\sim\!10^{9}$ ULP) sits five to
eight orders of magnitude \emph{above} the IEEE rounding floor. It is therefore
\emph{not} mere summation non-associativity --- there is a genuine, geometric
cause, isolated next.

\paragraph{Cause: the streamwise periodic seam is not bit-closed.}
The $\tfrac18$ dual telescopes \emph{exactly} in two of three directions: the
spanwise $x$ wrap is closed by the explicit remap $i{=}3\!\to\!N_x{-}5$ onto a
\emph{real} interior cell (and the in-plane area is $x$-independent, so all $256$
spanwise copies are bit-identical), and the wall-normal $k$ stencil is one-sided
at the walls with no periodic image required. The streamwise $y$ direction,
however, is closed only \emph{implicitly}, through the periodic ghost: the seam
node ($j{=}3$, Fr\"ohlich $I{=}0$) draws its left dual cell from the wrapped ghost
cell $j_{\mathrm{cell}}{=}2$, whereas $V_{\mathrm{disc}}$ counts the real last
cell $j_{\mathrm{cell}}{=}514$ (between nodes $I{=}512$ and the seam). These two
cells are translates by $L_y$ \emph{only if} the grid is bit-periodic. The
production Fr\"ohlich grid is not: nodes $I{=}0$ and $I{=}512$ differ from exact
periodic images by
\[
\max_k\bigl|y_{I=0}{+}L_y-y_{I=512}\bigr| = 3.285\times10^{-6},
\qquad \max_k|z_{I=0}-z_{I=512}| = 0 ,
\]
so the seam ghost area and the real last-cell area differ by up to
$2.62\times10^{-4}$ per wall-normal layer. Carrying that mismatch through the
$\tfrac12$ seam coefficient $\times\,(N_x{-}1)\,\Delta x$ predicts a closure
defect of $1.432306\times10^{-5}$, which reproduces the \emph{measured}
$\Sigma_w-V_{\mathrm{disc}}=1.432315\times10^{-5}$ to a relative
$6.5\times10^{-6}$. The seam thus accounts for essentially the entire gap.

\paragraph{Summation order (MPI non-associativity) is a sub-effect.}
Re-summing the same per-cell values in the real run's $jp{=}32$ grouping ($32$
contiguous streamwise rank-blocks of $16$, reduced in rank order to mimic
\texttt{MPI\_Reduce}) shifts each total only in its last bits:
$\Sigma_w$ moves by $28{,}313$ ULP ($4.02\times10^{-10}$ abs,
$3.52\times10^{-12}$ rel) and $V_{\mathrm{disc}}$ by $21{,}543$ ULP
($3.06\times10^{-10}$ abs, $2.68\times10^{-12}$ rel) --- the genuine
floating-point non-associativity caveat, of the order $n\varepsilon$ anticipated
in \eqref{eq:na-rounding}, and $\sim\!5$ orders smaller than the seam.

\paragraph{Closed-seam control: the Proposition holds to rounding.}
Forcing perfect periodic closure (node $I{=}0:=$ node $I{=}512-L_y$) and
recomputing \emph{both} sums on the closed grid collapses the discrepancy to the
rounding floor:
\[
|\Sigma_w-V_{\mathrm{disc}}|_{\text{closed}} = 9.27\times10^{-11},\quad
\text{rel}=8.11\times10^{-13}\ (<10^{-12}),\quad 6{,}522\ \text{ULP},
\]
i.e.\ the real-grid $1.25\times10^{-7}$ falls to $8.1\times10^{-13}$ once the seam
is bit-closed. This is the clean experimental statement of
Proposition~\ref{prop:pou}: $\sum_p w_p=\sum_c V_c$ is an \emph{exact-arithmetic}
identity that holds to IEEE-double rounding ($\sim\!10^{-13}$) --- \emph{not}
bitwise --- on a bit-closed periodic mesh, and is degraded only by the grid's own
seam-closure tolerance, not by the algorithm.

\begin{table}[t]
\centering
\small
\begin{tabular}{lccc}
\toprule
contribution to $\Sigma_w-V_{\mathrm{disc}}$ & abs & rel & ULP \\
\midrule
streamwise grid seam non-closure (geometry) & $1.43\times10^{-5}$ & $1.25\times10^{-7}$ & $1.0\times10^{9}$ \\
$jp{=}1$ vs $jp{=}32$ summation order (MPI)  & $\sim\!4\times10^{-10}$ & $\sim\!3\times10^{-12}$ & $\sim\!2$--$3\times10^{4}$ \\
IEEE rounding floor (closed-seam control)    & $9.3\times10^{-11}$ & $8.1\times10^{-13}$ & $6.5\times10^{3}$ \\
\bottomrule
\end{tabular}
\caption{Three-tier decomposition of the Method-0 partition-of-unity gap on the
production grid. The measured $1.25\times10^{-7}$ closure defect is
\emph{geometry}, dominated by the grid's streamwise seam ($\sim\!3\times10^{-6}$
non-closure); floating-point non-associativity (summation/MPI order) is the
$\sim\!10^{-12}$ sub-effect, and the bit-closed mesh recovers the
$\sim\!10^{-13}$ rounding floor.}
\label{tab:E4-tiers}
\end{table}

\paragraph{Comparison to the analytic hill volume.}
The continuum reference $V_{\mathrm{phys}}=L_x\bigl(L_yL_z-\int_0^{L_y}h\,\mathrm
dy\bigr)=1.143590335916\times10^{2}$ (Simpson, $N{=}10^4$, mirroring
\texttt{VerifyPhysicalDomainVolume3D}) is approached by both discrete volumes at
the $O(h^2)$ chord level: $V_{\mathrm{disc}}$ differs from $V_{\mathrm{phys}}$ by
$8.27\times10^{-7}$ (rel) and $\Sigma_w$ by $7.02\times10^{-7}$ --- the expected
geometric chord error of the curved hill, an order of magnitude larger than the
seam defect and three orders larger than the FP floor.

\paragraph{Method~1 (the default) departs by design.}
The shipped configuration is \texttt{CELL\_VOLUME\_METHOD}${=}1$. Reproducing the
real metric-term FD stencil (\texttt{ComputeMetricTerms\_Full}; six-order
Fornberg in $k$, six-order central in $j$, periodic-$j$ ghost) to obtain $J_{2D}$
on the production grid and applying the GL-3 / Lagrange-6 weight sum gives
$\Sigma_w^{\mathrm{M1}}=1.14359016044\times10^{2}$, which departs from
$V_{\mathrm{disc}}$ by $6.73\times10^{-7}$ (rel) --- a \emph{deliberate} $O(h^2)$
difference, not a rounding effect, because Method~1 integrates the curved
$J_{2D}$ rather than its straight chords (per-cell $|J\!-\!\text{Shoe}|/\text{Shoe}$
up to $1.78\times10^{-2}$, mean $9.21\times10^{-5}$; zero Shoelace fallbacks, the
smooth $J_{2D}{>}0$ never tripping the guard). Consistently with E2 and
\S\ref{sec:na-N1}, Method~1 lies \emph{closer} to the continuum volume:
$|\Sigma_w^{\mathrm{M1}}-V_{\mathrm{phys}}|/V_{\mathrm{phys}}=1.53\times10^{-7}$
versus Method~0's $8.27\times10^{-7}$.

\paragraph{Reconciliation and verdict.}
The audit confirms the precise status of the partition-of-unity identity on real
hardware. \emph{It is not bitwise} (\,$==$ is \textsc{false}, $\sim\!10^{9}$
ULP\,): $\Sigma_w$ equals the discrete main-field volume to a relative
$1.25\times10^{-7}$, dominated by the production grid's streamwise seam
non-closure ($\sim\!3\times10^{-6}$ in $y$) and \emph{not} by the
$\sim\!10^{-12}$ floating-point non-associativity of the $\tfrac18$-split
re-summation (the latter amplified, as a caveat, by \texttt{MPI\_Reduce} rank
order). Because $1.25\times10^{-7}$ (Method~0) --- and $6.73\times10^{-7}$ for
the default Method~1 --- both exceed the solver's documented $10^{-12}$ gate,
\texttt{VerifyPhysicalDomainVolume3D} emits its \texttt{WARNING}, not
\texttt{PASS}, on this grid; the warning is faithful, reporting a real
($O(h^2)$/seam-level) geometric reality, not a defect of the weight algebra. The
closed-seam control (rel $8.1\times10^{-13}{<}10^{-12}$,
Table~\ref{tab:E4-tiers}) shows Proposition~\ref{prop:pou} holds to rounding the
instant the mesh is bit-closed, exactly as the exact-arithmetic proof predicts.

% ===========================================================================
\section{Implementation: mapping the mathematics to source code}
\label{sec:impl}

This section maps every symbol of the volume-weighted mass-correction formulation
to the concrete C++/CUDA that realizes it. All references are to the Periodic-Hill
D3Q19 GILBM solver (\path{Edit13_2800ITBLBM}); line numbers are 1-indexed into the
cited files.

\subsection{Theory\texorpdfstring{$\leftrightarrow$}{<->}code correspondence}
\label{sec:impl:table}

\begin{table}[ht]
\centering
\small
\resizebox{\textwidth}{!}{%
\begin{tabular}{@{}lll@{}}
\toprule
Symbol / step & Realizing function & Location \\
\midrule
$V_{\text{cell}}=|\Delta x_i|\,A_{\text{quad}}$ (Shoelace) & \texttt{MassCorrectionCellVolume} & \path{evolution.h}:82--103 \\
$w_p=\tfrac18\sum_{c\ni p}V_{\text{cell}}$ (median-dual) & \texttt{InitializeMassCorrectionWeights} & \path{evolution.h}:105--172 \\
GL-3 nodes $\xi_a$, weights $W_a$ on $[0,1]$ & \texttt{GL3\_nodes}, \texttt{GL3\_weights} & \path{evolution.h}:180--189 \\
Degree-5 Lagrange basis $\ell_m(x)$ & \texttt{Lagrange6Weights} & \path{evolution.h}:191--204 \\
Centered/clamped 6-wide stencil start & \texttt{SelectStencilStart} & \path{evolution.h}:206--213 \\
$\tilde J_{2D}=\sum_{m,n}\ell_m\ell_n J_{2D}$ & \texttt{InterpolateJ2D\_Lagrange6} & \path{evolution.h}:215--231 \\
$A_{\text{jac}}\approx\sum_{a,b}W_aW_b\tilde J_{2D}$ & \texttt{ComputeJacobianMassCorrectionWeights} & \path{evolution.h}:233--362 \\
$V_{\text{mesh}}=\sum_c V_{\text{cell}}$ (discrete check) & \texttt{ComputeGlobalDiscreteShoelaceVolume3D} & \path{evolution.h}:364--383 \\
RelErr $<10^{-12}$ vs.\ discrete mesh & \texttt{VerifyPhysicalDomainVolume3D} & \path{evolution.h}:390--422 \\
$\sum_p w_p\rho_p$ (interior reduction) & \texttt{ReduceRhoWeightedSum\_Kernel} & \path{evolution.h}:47--80 \\
$\langle\rho\rangle_V=(\sum_p w_p\rho_p)/V$ & \texttt{ComputeVolumeWeightedRhoAverageRoot} & \path{evolution.h}:424--443 \\
$\delta\rho=1-\langle\rho\rangle_V$ + \texttt{MPI\_Bcast} & \texttt{UpdateVolumeWeightedMassCorrection} & \path{evolution.h}:445--453 \\
$\rho\!\leftarrow\!\rho+\delta\rho$, $f_0\!\leftarrow\!f_0+\delta\rho$ & fused collision kernel (apply site) & \path{gilbm/evolution\_gilbm/1.algorithm1.h}:407--411 \\
$V=\sum_p w_p$ via \texttt{MPI\_Allreduce} & \texttt{rho\_cv\_global\_volume} & \path{evolution.h}:149,333 \\
$J_{2D}=\det\partial(y,z)/\partial(\xi,\zeta)$ & \texttt{ComputeMetricTerms\_Full} & \path{gilbm/metric\_terms.h}:160--161 (\#incl.\ at \path{main.cu}:233) \\
ghost-row exchange of $J_{2D}$ (6-pt stencil) & \texttt{MPI\_Sendrecv} loop & \path{main.cu}:743--749 (list :741) \\
method selection + verification wiring & init block & \path{main.cu}:763--769 \\
compile-time method switch (default $1$) & \texttt{CELL\_VOLUME\_METHOD} & \path{variables.h}:342--347 \\
\bottomrule
\end{tabular}}
\caption{Each mathematical object maps to exactly one function/line range in the
source.}
\label{tab:corr}
\end{table}

\subsection{Key kernels (verbatim)}
\label{sec:impl:listings}

\noindent\emph{Note.} In the listings below the executable lines are reproduced
verbatim from the solver source; only the inline comments have been translated from
the original Traditional-Chinese for readability (e.g.\ the apply-site comment in
\path{gilbm/evolution\_gilbm/1.algorithm1.h}:407--411).

\subsubsection{Shoelace cell volume}
The planar quadrilateral area uses the surveyor (Gauss) formula on the four corners
$(y,z)$ at $(j,k),(j{+}1,k),(j{+}1,k{+}1),(j,k{+}1)$, multiplied by the spanwise
width.

\begin{lstlisting}[caption={Shoelace cell volume (\texttt{evolution.h}:82--103).}]
static inline double MassCorrectionCellVolume(const int i_cell, const int j_cell, const int k_cell)
{
    const double dx = x_h[i_cell + 1] - x_h[i_cell];

    const int jk00 = j_cell * NZ6 + k_cell;
    const int jk10 = (j_cell + 1) * NZ6 + k_cell;
    const int jk11 = (j_cell + 1) * NZ6 + (k_cell + 1);
    const int jk01 = j_cell * NZ6 + (k_cell + 1);

    const double y0 = y_2d_h[jk00], z0 = z_h[jk00];
    const double y1 = y_2d_h[jk10], z1 = z_h[jk10];
    const double y2 = y_2d_h[jk11], z2 = z_h[jk11];
    const double y3 = y_2d_h[jk01], z3 = z_h[jk01];

    const double area_yz = 0.5 * fabs(
          y0 * z1 - z0 * y1
        + y1 * z2 - z1 * y2
        + y2 * z3 - z2 * y3
        + y3 * z0 - z3 * y0);

    return fabs(dx) * area_yz;
}
\end{lstlisting}

\noindent\textbf{Annotation.} The four \texttt{jk}* indices fetch the corner
coordinates from the $y$--$z$ section arrays \texttt{y\_2d\_h}/\texttt{z\_h}.
\texttt{area\_yz} is the cyclic Shoelace sum
$\tfrac12\big|\sum y_a z_{a{+}1}-z_a y_{a{+}1}\big|$, the \emph{exact} area of the
straight-chord quadrilateral. The return value $|\Delta x_i|\cdot A_{\text{quad}}$
realizes $V_{\text{cell}}=|\Delta x_i|\,A_{\text{quad}}$. Because adjacent quads
share edges that cancel, $\sum_c V_{\text{cell}}$ telescopes to the exact
discrete-mesh polygon volume; this is precisely the quantity
\texttt{VerifyPhysicalDomainVolume3D} checks against (printing \texttt{PASS}) to $<10^{-12}$.

\subsubsection{GL-3 data, Lagrange basis, and \texorpdfstring{$J_{2D}$}{J2D} reconstruction}

\begin{lstlisting}[caption={GL-3 nodes/weights, degree-5 Lagrange basis, and tensor reconstruction (\texttt{evolution.h}:180--231).}]
static const double GL3_nodes[3] = {
    0.5 * (1.0 - 0.7745966692414834),   // (1 - sqrt(3/5))/2 ~ 0.1127
    0.5,
    0.5 * (1.0 + 0.7745966692414834)    // (1 + sqrt(3/5))/2 ~ 0.8873
};
static const double GL3_weights[3] = {
    5.0 / 18.0,    // ~ 0.27778
    8.0 / 18.0,    // ~ 0.44444
    5.0 / 18.0
};

static inline void Lagrange6Weights(double x, int start, double w[6])
{
    for (int m = 0; m < 6; m++) {
        double L = 1.0;
        double xm = (double)(start + m);
        for (int r = 0; r < 6; r++) {
            if (r != m) {
                double xr = (double)(start + r);
                L *= (x - xr) / (xm - xr);
            }
        }
        w[m] = L;
    }
}

static inline double InterpolateJ2D_Lagrange6(
    double xi_pos, double zeta_pos, int sj, int sk,
    const double *J_2D, int NZ6_local)
{
    double wj[6], wk[6];
    Lagrange6Weights(xi_pos,   sj, wj);
    Lagrange6Weights(zeta_pos, sk, wk);

    double result = 0.0;
    for (int m = 0; m < 6; m++) {
        double row_sum = 0.0;
        for (int n = 0; n < 6; n++)
            row_sum += wk[n] * J_2D[(sj + m) * NZ6_local + (sk + n)];
        result += wj[m] * row_sum;
    }
    return result;
}
\end{lstlisting}

\noindent\textbf{Annotation.} The constant \texttt{0.7745966692414834} is
$\sqrt{3/5}$, so \texttt{GL3\_nodes} are the three Gauss--Legendre abscissae mapped
onto $[0,1]$, $\{(1-\sqrt{3/5})/2,\ \tfrac12,\ (1+\sqrt{3/5})/2\}$, with weights
$\{\tfrac{5}{18},\tfrac{8}{18},\tfrac{5}{18}\}$ written exactly as
\texttt{5.0/18.0} etc. GL-3 is exact for polynomials of degree $\le 5$ per
direction. \texttt{Lagrange6Weights} builds the six cardinal degree-5 basis
functions $\ell_m(x)=\prod_{r\neq m}(x-x_r)/(x_m-x_r)$ on integer nodes
$\{start,\dots,start{+}5\}$. \texttt{InterpolateJ2D\_Lagrange6} forms the tensor
product $\tilde J_{2D}=\sum_{m}w_j[m]\sum_{n}w_k[n]\,J_{2D}[(s_j{+}m,s_k{+}n)]$.
Since the reconstruction is degree-5 and GL-3 integrates degree-5 exactly, the
interpolation does \emph{not} reduce the quadrature order.

\subsubsection{Jacobian-GL nodal weights (core triple loop)}

\begin{lstlisting}[caption={Median-dual assembly with GL-3 cells and Shoelace fallback (\texttt{evolution.h}:233--362, excerpt).}]
for (int j = 3; j < NYD6 - 4; j++) {
for (int k = 3; k < NZ6  - 3; k++) {
for (int i = 3; i < NX6  - 4; i++) {
    const int idx = j * NX6 * NZ6 + k * NX6 + i;

    const int i_cells[2] = { (i == 3) ? (NX6 - 5) : (i - 1), i };   // periodic x wrap
    const int j_cells[2] = { j - 1, j };
    int k_cells[2]; int nk_cells = 0;
    if (k > 3)       k_cells[nk_cells++] = k - 1;   // one-sided near wall
    if (k < NZ6 - 4) k_cells[nk_cells++] = k;

    double weight = 0.0;
    for (int jj = 0; jj < 2; jj++) {
    for (int ii = 0; ii < 2; ii++) {
    for (int kk = 0; kk < nk_cells; kk++) {
        const int jc = j_cells[jj];
        const int kc = k_cells[kk];
        const double dx = x_h[i_cells[ii] + 1] - x_h[i_cells[ii]];

        int sj, sk;
        bool sj_ok = SelectStencilStart(jc, j_lo_J, j_hi_J, &sj);
        bool sk_ok = SelectStencilStart(kc, k_lo_J, k_hi_J, &sk);

        double vol_jac; bool used_fallback = false;
        if (sj_ok && sk_ok) {
            double area_jac = 0.0;
            for (int a = 0; a < 3; a++) {
            for (int b = 0; b < 3; b++) {
                double xi_pos   = (double)jc + GL3_nodes[a];
                double zeta_pos = (double)kc + GL3_nodes[b];
                double J_val = InterpolateJ2D_Lagrange6(xi_pos, zeta_pos, sj, sk, J_2D, NZ6);
                if (!std::isfinite(J_val) || J_val <= 0.0) { used_fallback = true; break; }
                area_jac += GL3_weights[a] * GL3_weights[b] * J_val;
            }
            if (used_fallback) break;
            }
            vol_jac = fabs(dx) * area_jac;
        } else { used_fallback = true; }

        if (used_fallback) {
            vol_jac = MassCorrectionCellVolume(i_cells[ii], jc, kc);   // Shoelace fallback
            local_fallback_count++;
        }

        double vol_shoe = MassCorrectionCellVolume(i_cells[ii], jc, kc);
        if (vol_shoe > 0.0) {
            double rd = fabs(vol_jac - vol_shoe) / vol_shoe;
            if (rd > local_max_rel_diff) local_max_rel_diff = rd;
            local_sum_rel_diff += rd; local_cell_count++;
        }
        weight += 0.125 * vol_jac;
    }}}
    rho_cv_weight_h[idx] = weight;
    local_weight_sum += weight;
}}}
\end{lstlisting}

\noindent\textbf{Annotation.} The loop visits each interior unique node \texttt{idx}
and accumulates the median-dual weight $w_p=\tfrac18\sum_{c\ni p}V_{\text{cell}}$ via
the \texttt{0.125 * vol\_jac} contribution of up to eight surrounding cells
(\texttt{jj,ii,kk}). The periodic $x$ wrap is the ternary
\texttt{(i == 3) ? (NX6 - 5) : (i - 1)}; near-wall $k$ degrades to a one-sided
single cell (\texttt{nk\_cells} can be $1$). For each cell the inner $3\times3$ loop
forms $A_{\text{jac}}=\sum_{a,b}W_aW_b\,\tilde J_{2D}(jc{+}\xi_a,\ kc{+}\zeta_b)$ with
$\tilde J_{2D}$ from \texttt{InterpolateJ2D\_Lagrange6}, and the cell volume is
$|\Delta x|\cdot A_{\text{jac}}$. Any non-finite or non-positive reconstructed $J$
trips \texttt{used\_fallback}, which substitutes the exact Shoelace volume and
increments \texttt{local\_fallback\_count}. Independently, every cell records the
per-cell relative difference $|V_{\text{jac}}-V_{\text{shoe}}|/V_{\text{shoe}}$
feeding the \texttt{[MASS-CORR]} max/mean report.

\subsubsection{Average, correction, and the rest-population apply site}

\begin{lstlisting}[caption={Correction scalar, broadcast, and the rest-population apply site (\texttt{evolution.h}:445--453; \texttt{gilbm/evolution\_gilbm/1.algorithm1.h}:407--411).}]
// evolution.h:445-453  --- delta_rho = 1 - <rho>_V, broadcast to all ranks
static inline void UpdateVolumeWeightedMassCorrection()
{
    const double rho_avg = ComputeVolumeWeightedRhoAverageRoot();  // (sum w*rho)/V on root
    if (myid == 0) {
        rho_modify_h[0] = 1.0 - rho_avg;
    }
    MPI_Bcast(rho_modify_h, 1, MPI_DOUBLE, 0, MPI_COMM_WORLD);
    CHECK_CUDA( cudaMemcpy(rho_modify_d, rho_modify_h, sizeof(double), cudaMemcpyHostToDevice) );
}

// gilbm/evolution_gilbm/1.algorithm1.h:407-411  --- apply at every interior node, in the fused kernel
    // -- STEP 1.5: Macroscopic (mass correction) --
    if (i < NX6 - 4 && j < NYD6 - 4) {
        rho_stream += rho_modify[0];
        f_arr[0]   += rho_modify[0];   // correct f_arr[0] held in register
    }
\end{lstlisting}

\noindent\textbf{Annotation.} \texttt{rho\_modify\_h[0] = 1.0 - rho\_avg}
implements $\delta\rho=1-\langle\rho\rangle_V$ on the root rank; the
\texttt{MPI\_Bcast} distributes the single scalar so every rank applies an
identical shift. At the apply site the density moment receives
$\rho\!\leftarrow\!\rho+\delta\rho$ and the \emph{rest} population receives
$f_0\!\leftarrow\!f_0+\delta\rho$. Because $\mathbf e_0=(0,0,0)$, the injected
momentum is $\Delta(\rho u_\alpha)=\sum_q e_{q,\alpha}\Delta f_q=0$: mass is
injected, momentum is exactly preserved. The shift is spatially uniform, so the new
average is $\langle\rho\rangle_V+\delta\rho=1$ in a single application (the only
residual is floating-point rounding of the parallel reduction)---a discrete ``mass
governor'' with unit feedback gain.

\subsection{Implementation notes}
\label{sec:impl:notes}

\paragraph{Compile-time method switch and \texttt{\#error} guard.}
\path{variables.h}:342--347 defines the selector with a hard guard:
\begin{lstlisting}[caption={Method selector with compile-time guard (\texttt{variables.h}:342--347).}]
#ifndef CELL_VOLUME_METHOD
#define     CELL_VOLUME_METHOD  1
#endif
#if CELL_VOLUME_METHOD != 0 && CELL_VOLUME_METHOD != 1
#error "CELL_VOLUME_METHOD must be 0 (Shoelace) or 1 (Jacobian 3x3 GL)"
#endif
\end{lstlisting}
The default is \texttt{1} (Jacobian-GL). The entire
\texttt{ComputeJacobianMassCorrectionWeights} body is wrapped in
\texttt{\#if CELL\_VOLUME\_METHOD == 1} (\path{evolution.h}:233--362), so Method~0
builds compile only \texttt{InitializeMassCorrectionWeights}. \path{main.cu}:763--769
wires the choice: under Method~1 it calls
\texttt{ComputeJacobianMassCorrectionWeights(J\_2D\_h, \&shoelace\_vol)} then verifies
\emph{both} the Shoelace and Jacobian-GL volumes; otherwise it verifies the single
Shoelace volume.

\paragraph{Shoelace fallback path and diagnostics.}
Method~1 falls back to the exact Shoelace area whenever the 6-wide Lagrange stencil
is unavailable (\texttt{SelectStencilStart} returns \texttt{false} near a partition
edge: it needs \texttt{max\_start = hi - 5} with \texttt{max\_start - lo >= 0}) or
whenever a reconstructed $J_{2D}$ value is non-finite or $\le 0$. The root rank
prints the audit block (\path{evolution.h}:348--358): \texttt{Shoelace global
volume}, \texttt{Jacobian global volume}, $\Delta(\text{Jac-Shoe})/\text{Shoe}$,
\texttt{Per-cell max rel diff}, \texttt{Per-cell mean rel diff}, and
\texttt{Fallback to Shoelace = N cells}. \texttt{VerifyPhysicalDomainVolume3D}
additionally checks $\sum_p w_p$ against the telescoped discrete-mesh volume
\texttt{ComputeGlobalDiscreteShoelaceVolume3D}, prints \texttt{PASS} when
\texttt{rel\_err\_discrete < 1e-12} else a \texttt{WARNING}, and reports the
analytic $L_X(L_Y L_Z-\!\int_0^{L_Y}\!h\,\mathrm dy)$ Hill integral (Simpson,
$N=10000$) as a reference only.

\paragraph{MPI reductions.}
Total control volume $V=\sum_p w_p$ is a collective \texttt{MPI\_Allreduce}
(\path{evolution.h}:149 for Method~0, :333 for Method~1) into
\texttt{rho\_cv\_global\_volume} so every rank holds $V$. The weighted mass
$\sum_p w_p\rho_p$ is gathered by an \texttt{MPI\_Reduce} to root in
\texttt{ComputeVolumeWeightedRhoAverageRoot} (\path{evolution.h}:439); root forms
$\langle\rho\rangle_V$ and $\delta\rho$, then \texttt{MPI\_Bcast}
(\path{evolution.h}:451) distributes the scalar. The 6-point Lagrange stencil
reaches into neighbouring streamwise partitions, so the seven-array list at
\path{main.cu}:741 includes \texttt{J\_2D\_h} and the ghost-row
\texttt{MPI\_Sendrecv} loop at \path{main.cu}:743--749 exchanges
$\{$\texttt{xi\_y\_h, xi\_z\_h, zeta\_y\_h, zeta\_z\_h, y\_2d\_h, z\_h, J\_2D\_h}$\}$;
the Jacobian-weight code then allows $j$-stencils into the ghost rows
(\texttt{j\_lo\_J = 0}, \texttt{j\_hi\_J = NYD6-1}) while keeping $k$-stencils on
physical wall-normal nodes (\texttt{k\_lo\_J = 3}, \texttt{k\_hi\_J = NZ6-4}).

\paragraph{The \texorpdfstring{$f_0$}{f0} rest-population injection invariant.}
The reduction kernel \texttt{ReduceRhoWeightedSum\_Kernel}
(\path{evolution.h}:47--80) sums $\rho_p\,w_p$ over exactly the interior unique-node
box that the corrector later modifies ($i=3..\text{NX6}{-}5$,
$j=3..\text{NYD6}{-}5$, $k=3..\text{NZ6}{-}4$), guaranteeing the averaged and
corrected sets coincide. With D3Q19 ($N_Q=19$, $W_0=1/3$, $\mathbf e_0=\mathbf0$)
and the current \texttt{FORCE\_HERMITE\_ORDER = 2} (full Guo, whose first-order
density moment vanishes), the body force injects no net mass; the corrector's
$f_0\mathrel{+}=\delta\rho$ is the sole, momentum-neutral channel that drains the
accumulated compressibility drift back to $\langle\rho\rangle_V=1$.

% ===========================================================================
\section{Verification of the collision operator}
\label{sec:mrt}

For the volumetric corrector to be the \emph{sole} controlled mass source, the
collision operator and the body force must inject no spurious mass. We verified
this for the MRT collision on the conserved-moment basis of the D3Q19 transform
$M$, whose rows were checked numerically (not by annotation): row $k{=}0$ is the
all-ones density vector and rows $k{=}3,5,7$ equal the lattice-velocity components
$e_x,e_y,e_z$. The corresponding relaxation rates are pinned to zero,
$s_0{=}s_3{=}s_5{=}s_7{=}0$, so the precomputed projection $K=M^{-1}SM$ leaves these
moments untouched; we confirmed
$\max_{k\in\{0,3,5,7\}}\|M_k K\|_\infty = 1.66\times10^{-15}$, i.e.\ the collision
exactly conserves mass and momentum to machine precision. The Guo body force,
injected at second Hermite order
$F_i = W_i\,\mathcal{F}\,[\,3(c_{y}-v)+9(\mathbf{c}\!\cdot\!\mathbf{u})c_{y}\,]$ via
the split tables $(F_{\mathrm{proj}},F_{\mathrm{proj}}^{u,v,w})$ with the half-force
factor $(1-s_k/2){=}1$, satisfies $\sum_i F_i = \mathcal{O}(10^{-17})$ (zero net mass
source) while delivering exactly the prescribed streamwise momentum
$\sum_i F_i c_{y}=1$. Consequently the lattice collision contributes no spurious
mass, and the volumetric density corrector of \S\ref{sec:theory} remains the sole
controlled source of mass in the scheme. The full machine-checked record is
reproduced in Appendix~\ref{app:mrt}.

% ===========================================================================
\section{Discussion and limitations}
\label{sec:discussion}

\subsection{Discrete self-consistency, not field fidelity, is the correctness criterion}
\label{sec:self-consistency}

The seam audit of \S\ref{sec:E4} invites a question of principle that we now
settle: against \emph{what} reference is the volume-weighted corrector to be
judged correct? The operative criterion is \emph{internal self-consistency on the
solver's own discrete grid}---not agreement with the true straight-boundary
physical domain, and not even exact partition-of-unity closure. We make this
precise.

\paragraph{The corrector is unconditionally self-consistent.}
The governor of \S\ref{sec:T4} forms $\langle\rho\rangle_V$ from
\eqref{eq:rho-mean} using the \emph{same} nodal weights $w_p$ in numerator and
denominator. Consequently, for any spatially uniform field $\rho_p\equiv c$,
\begin{equation}
\rho_p\equiv c
\quad\Longrightarrow\quad
\langle\rho\rangle_V
\;=\;\frac{\sum_p \rho_p\,w_p}{\sum_p w_p}
\;=\;\frac{c\,\sum_p w_p}{\sum_p w_p}
\;=\;c ,
\label{eq:disc-selfconsistency}
\end{equation}
to floating-point rounding and \emph{independently of the actual values of the
weights} $w_p$. In particular $\rho_p\equiv 1$ yields $\langle\rho\rangle_V=1$ and
hence $\delta\rho=1-\langle\rho\rangle_V=0$ by \eqref{eq:deltarho}: the corrector
injects nothing into a uniform field, and the one-shot map \eqref{eq:one-shot} has
already reached its fixed point \eqref{eq:fixed-point}. Crucially, this identity
holds \emph{unconditionally}---regardless of whether $\Sigma_w=\sum_p w_p$ equals
the discrete cell sum $\sum_c V_c$, and regardless of whether either equals the
true physical volume $V_{\mathrm{phys}}$. It is a property of the \emph{algebraic
form} of \eqref{eq:rho-mean}, not of any geometric accuracy of $w_p$.

\paragraph{Three distinct notions of ``volume agreement'' must not be conflated.}
The discussion of \S\ref{sec:T3-pou} and \S\ref{sec:E4} touches three logically
independent statements, only the first of which bears on correctness:
\begin{enumerate}
\item[(i)] \emph{Internal self-consistency}: numerator and denominator of
  \eqref{eq:rho-mean} share the weights $w_p$, so a uniform field is reproduced
  exactly and no spurious correction is injected. This is \eqref{eq:disc-selfconsistency};
  it holds \emph{unconditionally}, for \emph{any} positive weights.
\item[(ii)] \emph{Partition-of-unity closure}:
  $\Sigma_w=\sum_p w_p$ equals $\sum_c V_c=V_{\mathrm{mesh}}$ on the solver's own
  grid---Proposition~\ref{prop:pou}, equation~\eqref{eq:pou}. This is an
  exact-arithmetic identity that holds to IEEE-double rounding \emph{if and only
  if} the grid's periodic seam is bit-closed; it is exactly the quantity
  \texttt{VerifyPhysicalDomainVolume3D} measures, and exactly what the seam audit
  \eqref{eq:E4-metrics} probes.
\item[(iii)] \emph{Field fidelity}: $\Sigma_w$ equals $V_{\mathrm{phys}}$, the
  volume of the continuum straight-boundary periodic-hill domain. This is
  \emph{never} exact---the discrete boundary is a curve, and the
  interior outlet rows land at, e.g., $y\approx 8.9999967$ rather than $y=L_y$
  (the read-time normalization pins the anchor row exactly). The discrepancy is the
  $O(h^2)$ chord error of the curved hill plus the single-precision streamwise
  seam non-closure.
\end{enumerate}

\paragraph{Only (i) is required for correctness.}
Property~(i) is \emph{necessary and sufficient} for the corrector to do its job:
hold the volume-weighted mean density at unity without injecting a spurious
source. Property~(ii) is a desirable consistency and diagnostic property---it
certifies that $\Sigma_w$ is a faithful discrete mesh volume---but the corrector
is self-consistent and mass-conserving whether or not (ii) holds exactly, because
\eqref{eq:disc-selfconsistency} never invokes $\sum_c V_c$. Property~(iii) is a
statement of \emph{geometric accuracy}---how faithfully the discrete domain
approximates the continuum---and therefore bears on \emph{solution} accuracy, not
on mass-conservation \emph{correctness}. Mass is conserved on the grid the solver
is actually given (curved boundary, single-precision provenance and all): the
corrector holds $\langle\rho\rangle_V$ at unity with respect to its \emph{own}
control volumes $w_p$, which is the only mass that the discrete dynamics either
sees or can conserve.

\paragraph{The seam discrepancy is a geometric diagnostic, not a defect.}
It follows that the \texttt{VOL-CHECK} discrepancy reported in the E4 audit
($\Sigma_w$ versus $\sum_c V_c$, relative $\sim\!10^{-7}$,
\eqref{eq:E4-metrics}) is a \emph{geometric diagnostic of streamwise periodic
closure}---it flags that the production grid's outlet is not bit-closed onto its
inlet---and is \emph{not} a mass-conservation defect. The corrector never reads
$\sum_c V_c$; it reads $\Sigma_w$ in both numerator and denominator, so the seam
gap cancels identically in \eqref{eq:rho-mean} and cannot bias $\delta\rho$.
Closing the seam---E4's control, node $I{=}0:=$ node $I{=}512-L_y$, i.e.\ making
the two boundary edges exact periodic images using the grid's own
coordinates---makes (ii) exact: it collapses the discrepancy to the rounding floor
(\textsc{pass}, \S\ref{sec:E4}, Table~\ref{tab:E4-tiers}) and additionally
renders the grid bit-periodic. But this is a cleanliness and diagnostic
refinement, \emph{not} a correctness fix: the corrector was already
self-consistent and mass-conserving before closure. The grid is not ``wrong''---it
is a faithful $O(h^2)$/seam-level discrete approximation of the continuum hill---and
the corrector's correctness is established by \eqref{eq:disc-selfconsistency}
alone, leaving the seam a question of grid closure and geometry only.

\paragraph{Why volume weighting is not optional.}
The central message is structural: on a stretched body-fitted mesh, the only
mesh-invariant target for a density governor is the volume-weighted mean
$\langle\rho\rangle_V$. Regulating the unit-weight sum $\widetilde M$ of
\eqref{eq:Mtilde-wrong} is not a milder approximation but a qualitatively wrong
control objective---it pins a metric-biased functional and feeds a spurious,
mesh-dependent momentum source. The median-dual weights make the correct objective
computable at the cost of one precomputed weight array and one extra reduction per
step.

\paragraph{Cost.}
The corrector adds a single global reduction ($\sum_p w_p\rho_p$), a scalar
broadcast, and a per-node add fused into the existing collision kernel. The weight
array $w_p$ and the volume $V$ are built once at start-up; the GL-3/Lagrange
machinery is therefore amortised entirely into initialisation and does not enter
the time-stepping cost.

\paragraph{Limitations.}
(i) The correction is \emph{uniform}: it removes the global mean drift but does not
redistribute mass spatially, so a genuinely spatially-structured mass error (e.g.\
from a defective boundary condition) would not be cured---only its global mean. (ii)
The one-shot exactness \eqref{eq:na-oneshot} holds for the volume-weighted mean; the
residual is the non-associative parallel-reduction rounding of
\eqref{eq:na-rounding}, which makes $\delta\rho$ reproducible only at fixed rank
count and fixed reduction tree. (iii) Method~1's sixth-order accuracy presumes a
\emph{smooth} metric $J_{2D}$; near a geometric singularity the degree-5 Lagrange
reconstruction loses order and the code (correctly) falls back to the exact
straight-chord Shoelace area, recovering the discrete-mesh volume but reintroducing
the $O(h^2)$ geometric bias there. (iv) The whole construction exploits the block
factorisation \eqref{eq:jacobian-block} of the extruded geometry; a fully
three-dimensional curvilinear map would require a $3\times3$ Jacobian quadrature
rather than the present area$\times$chord separation, though the median-dual and
corrector theory carry over unchanged. (v) The host-only verification harness of
\S\ref{sec:experiments} reproduces the solver kernels (the GL-3 nodes/weights, the
degree-5 Lagrange basis, the centred/clamped stencil selector, the tensor
reconstruction, and the Shoelace area) as \emph{hand-copied snapshots}, because the
originals are CUDA/MPI-coupled and cannot be linked into a standalone program; the
harness must therefore be kept in manual sync with \path{evolution.h} and will
otherwise silently validate stale code---a known limitation that a shared header
would remove, at the cost of editing the solver source.

\paragraph{Generality.}
Nothing in \S\ref{sec:theory}--\S\ref{sec:weights} is specific to D3Q19 beyond the
existence of a zero-velocity rest population to carry the momentum-silent shift; any
lattice with $\mathbf e_0=\mathbf0$ admits the identical corrector, and any
body-fitted/finite-volume discretisation with a partition-of-unity nodal volume
admits the identical volume-weighted mean.

% ===========================================================================
\section{Conclusion}
\label{sec:conclusion}

We have formulated, analysed, verified, and documented a Jacobian-weighted
volumetric density corrector for mass conservation in curvilinear lattice-Boltzmann
flows. The theory establishes that the physical mass $M=\int_\Omega\rho\,dV$, and
hence the volume-weighted mean $\langle\rho\rangle_V$, is the only mesh-invariant
target, and proves that a uniform shift into the rest population $f_0$ restores
$\langle\rho\rangle_V=1$ in a single, momentum-preserving step. The median-dual
control volumes satisfy an exact partition-of-unity identity, and the cell areas
that build them are integrated either exactly (straight-chord Shoelace) or to sixth
order (Jacobian $3\times3$ Gauss--Legendre of a degree-5 Lagrange reconstruction of
the in-plane metric). An independent host-only harness confirms GL-3 degree-5
exactness, sixth-order convergence of the Jacobian method against a textbook
$O(h^2)$ for the chorded method, and algebraically exact, momentum-preserving mass
restoration to the reduction-rounding floor ($\sim\!10^{-16}$ in the small
verification harness, $\sim\!10^{-12}$ at production scale). A complete
theory$\leftrightarrow$code map ties
every symbol to the D3Q19 GILBM source, and a machine-checked record establishes
that the MRT collision and Guo forcing inject no spurious mass---leaving the
volumetric corrector as the sole, controlled mass channel. The result is a mass
governor that is mesh-invariant by construction, exact in one step in exact
arithmetic, and limited in practice only by the well-conditioned, tree-summed
parallel reduction.

% ===========================================================================
\appendix
\section{MRT collision verification record}
\label{app:mrt}

The following is the verbatim machine-checked verification record for the MRT
collision operator and Guo forcing used by the solver (mixed
English/Traditional-Chinese audit log, reproduced as run).

\begin{CJK}{UTF8}{bsmi}
\small
\begin{verbatim}
=== MRT/LBM Collision 啟動驗證 ===
結論: 正確

Conserved Moment Indices (由 M 矩陣數值證明，非靠變數名/註解):
  k=0  density  (s=0, M row = {1,1,1,...,1} 全為 1)
  k=3  rho*u_x  (s=0, M row == e_x = {0,1,-1,0,0,0,0,1,-1,1,-1,1,-1,1,-1,0,0,0,0})
  k=5  rho*u_y  (s=0, M row == e_y = {0,0,0,1,-1,0,0,1,1,-1,-1,0,0,0,0,1,-1,1,-1})  <- streamwise
  k=7  rho*u_z  (s=0, M row == e_z = {0,0,0,0,0,1,-1,0,0,0,0,1,1,-1,-1,1,1,-1,-1})

檢查結果:
  [v] Relaxation rates: conserved s[0]=s[3]=s[5]=s[7]=0.0
  [v] Collision: conserved moments 不被 relax (projection f_out=f_B-K(f-feq))
  [v] K 矩陣: max|M_conserved.K| = 1.66e-15 < 1e-12 (k=0,3,5,7 四列全部 ~0)
  [v] Forcing: density 矩=0 (sum_i F0=0)，y-momentum 矩=1，(1-s_k/2)=1 因 s_k=0
  [v] Force 未被覆蓋: f_out[a]=f_B[a]-relax 之後 += dt*Force0*(Fproj+u*Fu+v*Fv+w*Fw)
  [v] 註解正確性: u/v/w = spanwise/streamwise/wall-normal 標註正確

FORCE_HERMITE_ORDER=2 驗證:
  二階完整 Guo: F_i = w_q*Force*[3(c_y-v) + 9(c.u)c_y]
  sum_i F_i = 1.4e-17 ~ 0 (質量守恆)，sum_i F_i*c_y = 1.0 (y-動量守恆)
  Host 拆成 4 基底 F0=w*3c_y, Fu=w*9c_x c_y, Fv=w*(9c_y^2-3), Fw=w*9c_z c_y
  kernel 線性重組 Fproj + u*Fu + v*Fv + w*Fw 等價於完整二階 (容差 1e-12)
  一階若啟用為 w_q*3*c_y*Force（無 -v 項）

相關檔案:
  MRT matrix:     MRT_Matrix.h:26 (M), :48 (Mi), :4 (Relaxation s)
  Relaxation:     mrt_projection_host.h:20, variables.h:118 (USE_MRT), :331 (FORCE_HERMITE_ORDER=2)
  Collision:      gilbm/evolution_gilbm/0.collision.h:38-78, :65 (K), :69-73 (Guo force)
  Forcing build:  mrt_projection_host.h:127-170 (K+Fproj), :103-125, :264-416 (Verify)
  Constants:      gilbm/evolution_gilbm/0.shared_code.h:48-52 (__constant__ K, Fproj*)

Bug (若有): 無。質量與動量守恆矩、conserved s=0、K 投影零化、force 注入與半力係數、
            二階 Guo 拆解全部通過 1e-12 容差。

驗證細節: 以 numpy 重建 M、解析 Mi (Mi.M 恆等誤差 1.35e-15)、s=diag(0,1.19,...)，
          組 K=Mi.S.M。M[k].K 對 k in {0,3,5,7} 最大 1.66e-15。二階直接 Guo 力
          sum=1.4e-17、sum.c_y=1，確認 split 基底重組無誤。
\end{verbatim}
\end{CJK}

% ===========================================================================
\begin{thebibliography}{99}

\bibitem{Qian1992}
Y.~H. Qian, D.~d'Humi\`eres, and P.~Lallemand,
\newblock Lattice BGK models for Navier--Stokes equation,
\newblock \emph{Europhysics Letters} \textbf{17}(6), 479--484, 1992.

\bibitem{dHumieres2002}
D.~d'Humi\`eres, I.~Ginzburg, M.~Krafczyk, P.~Lallemand, and L.-S. Luo,
\newblock Multiple-relaxation-time lattice Boltzmann models in three dimensions,
\newblock \emph{Philosophical Transactions of the Royal Society A} \textbf{360}(1792),
360--451, 2002.

\bibitem{GuoZhengShi2002}
Z.~Guo, C.~Zheng, and B.~Shi,
\newblock Discrete lattice effects on the forcing term in the lattice Boltzmann method,
\newblock \emph{Physical Review E} \textbf{65}(4), 046308, 2002.

\bibitem{Kruger2017}
T.~Kr\"uger, H.~Kusumaatmaja, A.~Kuzmin, O.~Shardt, G.~Silva, and E.~M. Viggen,
\newblock \emph{The Lattice Boltzmann Method: Principles and Practice},
\newblock Springer, Cham, 2017.

\bibitem{He1997}
X.~He and L.-S. Luo,
\newblock Theory of the lattice Boltzmann method: From the Boltzmann equation to the
lattice Boltzmann equation,
\newblock \emph{Physical Review E} \textbf{56}(6), 6811--6817, 1997.

\bibitem{FilippovaHanel1998}
O.~Filippova and D.~H\"anel,
\newblock Grid refinement for lattice-BGK models,
\newblock \emph{Journal of Computational Physics} \textbf{147}(1), 219--228, 1998.

\bibitem{StoerBulirsch2002}
J.~Stoer and R.~Bulirsch,
\newblock \emph{Introduction to Numerical Analysis}, 3rd ed.,
\newblock Springer, New York, 2002.

\bibitem{DavisRabinowitz1984}
P.~J. Davis and P.~Rabinowitz,
\newblock \emph{Methods of Numerical Integration}, 2nd ed.,
\newblock Academic Press, Orlando, 1984.

\bibitem{Barth1991}
T.~J. Barth,
\newblock Aspects of unstructured grids and finite-volume solvers for the Euler and
Navier--Stokes equations,
\newblock \emph{AGARD/VKI Lecture Series} 1991-05 (von Karman Institute), 1991.

\bibitem{Frohlich2005}
J.~Fr\"ohlich, C.~P. Mellen, W.~Rodi, L.~Temmerman, and M.~A. Leschziner,
\newblock Highly resolved large-eddy simulation of separated flow in a channel with
streamwise periodic constrictions,
\newblock \emph{Journal of Fluid Mechanics} \textbf{526}, 19--66, 2005.

\bibitem{Breuer2009}
M.~Breuer, N.~Peller, C.~Rapp, and M.~Manhart,
\newblock Flow over periodic hills---Numerical and experimental study in a wide range
of Reynolds numbers,
\newblock \emph{Computers \& Fluids} \textbf{38}(2), 433--457, 2009.

\bibitem{Succi2001}
S.~Succi,
\newblock \emph{The Lattice Boltzmann Equation for Fluid Dynamics and Beyond},
\newblock Oxford University Press, Oxford, 2001.

\end{thebibliography}

\end{document}
